% framed knot への refined large N 双対の拡張と，トーラス結び目を超えた正値性の検証（日本語の準備稿）
% 構成：docs/framed_knot/paper/outline.md，主張の格付け：docs/framed_knot/paper/claims.md（各主張の直前に % [ID] で対応を書く），記号：docs/framed_knot/paper/notation.md
% コンパイル：cd docs/framed_knot/paper/ja && latexmk -lualatex main.tex
\documentclass[11pt,a4paper]{ltjsarticle}
\usepackage{amsmath,amssymb,amsthm}
\usepackage{booktabs,array,longtable,graphicx}
\usepackage[margin=25mm]{geometry}
\usepackage{youngtab}
\usepackage{tcolorbox}
% References: biblatex (numeric) so that arXiv numbers are printed; shorthand fields of the INSPIRE entries are not used as labels.
\usepackage[backend=bibtex,style=numeric,sorting=nyt,maxbibnames=10,doi=false,url=false,isbn=false,eprint=true]{biblatex}
\addbibresource{../references.bib}
\DeclareFieldInputHandler{shorthand}{\def\NewValue{}}
\AtEveryBibitem{\iffieldundef{eprint}{}{\clearfield{note}}}
\usepackage[hidelinks]{hyperref}

\newtheorem{theorem}{定理}[section]
\newtheorem{proposition}[theorem]{命題}
\newtheorem{lemma}[theorem]{補題}
\newtheorem{corollary}[theorem]{系}
\newtheorem{conjecture}[theorem]{予想}
\theoremstyle{definition}
\newtheorem{definition}[theorem]{定義}
\newtheorem{hypothesis}[theorem]{仮説}
\newtheorem{compresult}[theorem]{計算結果}
\theoremstyle{remark}
\newtheorem{remark}[theorem]{注意}

% ヤング図は KN と同じく youngtab の \yng(2,1) で書く
\Yboxdim4pt
\newcommand{\wh}{\widehat}
\newcommand{\U}{\mathrm{U}}
\newcommand{\SU}{\mathrm{SU}}
\newcommand{\rCS}{\textrm{rCS}}
\newcommand{\ba}{\mathbf{a}}
\newcommand{\bq}{\mathbf{q}}
\newcommand{\bt}{\mathbf{t}}
\newcommand{\hN}{\widehat{N}}
\newcommand{\hf}{\widehat{f}}
\newcommand{\hbN}{\widehat{\mathbf{N}}}
\newcommand{\Exp}{\mathrm{Exp}}

\title{framed knot への refined large $N$ 双対の拡張と\\トーラス結び目を超えた正値性の検証\\
\large（日本語の準備稿）}
\author{亀山 昌也（Masaya Kameyama）}
\date{2026年10月10日（準備稿）}

\begin{document}
\maketitle
\begin{abstract}
Kameyama--Nawata~\cite{Kameyama17}（以下 KN）は，refined Chern--Simons 不変量を resolved conifold 上の refined BPS 不変量に書き換え，
それらが非負整数になるという正値性予想を，framing 0 のトーラス結び目・トーラス絡み目と，いくつかの thin な結び目で確かめた．
本論文では，この枠組みを framing をもつ結び目に拡張し，KN が用いたデータを二つの方向に超えて，
refined LMOV 不変量の spin 表示 $\hbN_{\rho;i,j,k}$（genus 展開の前）と genus 表示 $\hbN_{\rho,g,\beta,F}$（genus 展開の後）の正整数性を検証する．

まず，framing $x$ の不変量を $\overline{\rCS}^{(x)}_\lambda=\theta_\lambda^x\,\overline{\rCS}_\lambda$ で定める．
$\theta_\lambda$ は Taki の refined topological vertex の framing の因子（の逆数）に一致し，$t=q$ で Mariño--Vafa の framing の規則に戻る．
$\bt^{-|\rho|x}\hf^{(x)}_\rho$ は $x$ の指数和であり，KN の枠組みの構造的な性質（転置対称性，Laurent 性，genus 展開が閉じる条件）は $x$ によらない．
Mariño--Vafa の符号 $(-1)^{|\lambda|x}$ は，Laurent 性を保つために必要であることを示す．

次に，結び目と framing の両方の方向で検証する．
(i) ホモロジーを入力とする限り，任意の結び目・任意の framing で基本表現 $\rho=\yng(1)$ の正値性が成り立つ（命題）．11 交点以下の全素結び目 801 個で確かめた．
(ii) 2 箱では，Laurent 性の必要条件 $\dim H^{\yng(2)}=(\dim H^{\yng(1)})^2$ を導き，予想式を入力として，ツイスト結び目の全ての撚りの数での正値性を示す．
また 8 個の結び目では，全ての整数の framing で正値性を示す．
(iii) 3 箱では，トーラス結び目 7 個と代数的な iterated torus 結び目 7 個（3 本ケーブルと 2 段の反復を含む）について，
定義どおりの refined 不変量から，$|\rho|\le 3$ の全ての $\rho$ で spin 表示と genus 表示の両方で正であることを計算で確かめた．
最大の例は約 183 万項である．
さらに framing を変えても正値性は保たれる．代数的な iterated torus 結び目とトーラス結び目で $|x|\le10$（$K10$ の 3 箱は $|x|\le2$）まで，計 1821 件を確かめた．
(iv) 4 箱以上では，任意の箱数で $\hf_\rho$ を計算する手順（KN の行列 $M(t)$ の対称群の指標による対角化）を用い，トーラス結び目 8 個で $|\rho|=4$ から最大 9（$T_{2,3}$）まで，計 340 個の $(K,\rho)$ で正であることを計算で確かめた．
$\mathbf{4_1}$ の 4 箱では，未知の $\yng(3,1)$ に条件を課すと，正値性を満たす有理数の解が存在する（整数解は未決）．
また，KN の $\mathbf{4_1}$ の $\yng(2,1)$ の候補が Anokhina らの $Z$ 因子展開の係数を一つ取り替えたものであることを示す．
これらに基づき，framed knot の refined LMOV 不変量の正値性を予想として提出する．
\end{abstract}

\tableofcontents

\section{序論}\label{sec:intro}

\subsection{refined large $N$ 双対と KN の正値性予想}
$U(N)$ Chern--Simons 理論の large $N$ 極限は resolved conifold $X=\mathcal{O}(-1)\oplus\mathcal{O}(-1)\to\mathbb{P}^1$ 上の位相的弦理論と等価であり，
結び目 $K$ を加えると，色付き HOMFLY-PT 多項式は $X$ の中の Lagrangian $L_K$ に端をもつ BPS 状態の数（LMOV 不変量）で書ける
\cite{Ooguri99,Labastida01}．
Kameyama--Nawata~\cite{Kameyama17} は，この描像を refined Chern--Simons 理論~\cite{Aganagic11} に広げた．
M 理論の refined index と Schwinger の計算から自由エネルギーの形を定め，refined CS 不変量 $\overline{\rCS}_\lambda(K;a,q,t)$ から
refined reformulated invariants $f^q_\mu$, $f^{\bar t}_\mu$ と $\hf_\rho$ を作る式を与えた．
そのうえで，$\hf_\rho$ を太字の変数 $(\ba,\bq,\bt)$ で書いたときの係数 $\hbN_{\rho;i,j,k}$（本論文の spin 表示）と，
$G=(\bq\bt-\bq^{-1}\bt^{-1})(\bq-\bq^{-1})$ の冪に展開（M2-brane の genus 展開）した後の係数 $\hbN_{\rho,g,\beta,F}$（本論文の genus 表示．後述の予想の下で refined BPS 不変量 $\hN_{\rho,g,\beta,J_r}$ に一致する）が，非負整数になることを予想した．
物理的には，この正値性は追加の $U(1)_R$ 対称性が BPS 状態に自明に作用することから期待される．

Kameyama--Nawata~\cite{Kameyama17} がこの予想を確かめたのは，framing 0 の次のデータである（以下「KN のデータ」と呼ぶ）．
トーラス結び目・トーラス絡み目の $|\lambda|\le 3$ の表現，
Gukov--Nawata--Saberi--Stošić--Sułkowski（GNSSS）~\cite{Gukov15} の $(2s-1,1,2t-1)$-pretzel 結び目（ツイスト結び目を含む）と $\mathbf{6_2},\mathbf{6_3}$ の $\yng(1),\yng(2),\yng(1,1)$（色付き HOMFLY-PT ホモロジーの Poincaré 多項式を入力とする），
および KN が構成した $\mathbf{4_1}$ の $\yng(2,1)$ の候補である．
KN はさらに次の二点を指摘した．
非長方形の表現の refined CS 不変量は負の係数をもつが，genus 表示は非負になること．
thick な（HOMFLY-PT ホモロジーが一つの $\delta$ 次数に集中しない）結び目 $\mathbf{9_{42}}$ は形 \eqref{eq:hat-f-2} に乗らず，それは thick なホモロジーが指数増大性をもたないこと（Gukov--Stošić~\cite{Gukov11}）と関係すること．
本論文では，KN のデータの再現は規約を定めるためだけに用い，KN で観察された現象は成果としては扱わない．

\subsection{本論文の内容}
本論文の目的は，KN の枠組みを framing をもつ結び目に拡張し，KN のデータを二つの方向に超えて正整数性を検証することである．
一つは結び目の方向で，トーラス結び目と thin な結び目から，代数的な iterated torus 結び目，結び目表の全ての結び目，ツイスト結び目・ダブルツイスト結び目の族，thick な結び目へ広げる．
もう一つは framing の方向で，$x=0$ から，有限の範囲の $x$，さらに全ての整数 $x$ へ広げる．

\paragraph{framed knot への拡張（\S\ref{sec:framing}）．}
framing $x$ の不変量を $\overline{\rCS}^{(x)}_\lambda=\theta_\lambda^x\,\overline{\rCS}_\lambda$，
\[
\theta_\lambda=\bigl(-(q/t)^{1/2}\bigr)^{|\lambda|}\,q^{n(\lambda')}\,t^{-n(\lambda)}
\]
で定める．この因子は Taki~\cite{Taki07} の refined topological vertex の framing の因子の逆数にちょうど一致し，
$t=q$ では Mariño--Vafa~\cite{Marino01} の framing の規則 $W_R\mapsto(-1)^{\ell p}q^{p\kappa_R/2}W_R$（$\ell=|R|$）に $x=p$ で戻る．
太字の変数では $\theta_\lambda$ は符号のない単項式 $\bt^{|\lambda|+2n(\lambda')}\bq^{2(n(\lambda')-n(\lambda))}$ になり，
framing はホモロジーの次数のずらしとして働く．
主な構造的な結果は次のとおりである．
\begin{itemize}
\item $\bt^{-|\rho|x}\hf^{(x)}_\rho$ は $x$ の指数和である（命題~\ref{prop:expsum}）．
\item 転置対称性，Laurent 性，genus 展開が閉じる条件は $x$ によらない（命題~\ref{prop:x-indep}）．
\item Mariño--Vafa の符号 $(-1)^{|\lambda|x}$ を外すと，奇数の $x$ で 2 箱の $\hf$ が Laurent でなくなる（命題~\ref{prop:sign}）．
\item KN の単項式の正規化のうち表現によって変わる部分は，ちょうど Seifert framing の分の framing である（命題~\ref{prop:N2-framing}）．
\end{itemize}

\paragraph{検証（\S\ref{sec:one}--\S\ref{sec:higher}）．}
入力の確からしさは結び目のクラスによって異なる（表~\ref{tab:overview}）．
\begin{itemize}
\item 基本表現 $\rho=\yng(1)$：入力が非負係数で Dunfield--Gukov--Rasmussen（DGR）~\cite{Dunfield05} の対称性 \eqref{eq:sigma} をもてば，genus 展開は整数で閉じ，前後とも非負である（命題~\ref{prop:one}）．
DGR の対称性は Gorsky--Hogancamp--Mellit~\cite{Gorsky21} の定理なので，ホモロジーを入力とする限り，任意の結び目・任意の framing で $\yng(1)$ は正である．
\item 2 箱：KN が指摘した $\mathbf{9_{42}}$ と指数増大性の関係を，Laurent 性の必要条件 $\dim H^{\yng(2)}=(\dim H^{\yng(1)})^2$ として証明する（命題~\ref{prop:two-necessary}）．
KN が観察したツイスト結び目とダブルツイスト（pretzel）結び目の族について，Kameyama--Nawata--Tao--Zhang（KNTZ）~\cite{Kameyama19} の refined の式（GNSSS の予想式と一致する）を入力として，
ツイスト結び目の全ての撚りの数で 2 箱が正であることを証明する（定理~\ref{thm:twist}．ダブルツイスト結び目は定理~\ref{thm:pretzel}）．
また，8 個の結び目では全ての整数の framing で 2 箱が正である（定理~\ref{thm:two-all-x}）．
\item 3 箱：トーラス結び目 7 個と，代数的な iterated torus 結び目 7 個で，正規化 (N1)(N2$'$)（\S\ref{sec:normalization}）の下で，$|\rho|\le3$ の全ての $\rho$ について spin 表示と genus 表示の両方で正である（計算結果~\ref{cr:three}）．
後者は Cherednik--Danilenko~\cite{Cherednik14} の DAHA による定義どおりの refined 不変量を入力とし，
3 本ケーブル $K7$ と 2 段の反復 $K10$（3 箱で約 183 万項）を含む（結び目の名前は表~\ref{tab:cables}）．
KN がトーラス結び目で観察したとおり，入力の $\rCS_{\yng(2,1)}$ 自体は負の係数をもつが genus 表示は正になる．本論文はこの現象を代数的な iterated torus 結び目に広げる．
\item 4 箱以上：KN の行列 $M(t)$ を対称群の指標で対角化し（補題~\ref{lem:M-diag}），任意の箱数で $\hf_\rho$ を計算する．
トーラス結び目 8 個で，$|\rho|=4$ から最大 9 まで（計 340 個の $(K,\rho)$）正である（計算結果~\ref{cr:higher}）．
$\mathbf{4_1}$ の 4 箱は，未知の $\yng(3,1)$ に条件を課した弱い判定にとどまる（計算結果~\ref{cr:fig8-four}）．
\end{itemize}
これらに基づき，framed knot の refined LMOV 不変量の正値性を予想~\ref{conj:A} として，
代数的な iterated torus 結び目の正値性を予想~\ref{conj:B} として提出する．
双曲結び目の 3 箱には，定義に基づく $\yng(2,1)$ の refined 不変量が無い．構造的な条件では $\yng(2,1)$ が一意に決まらないことを \S\ref{sec:21} で述べる．
KN の $\mathbf{4_1}$ の $\yng(2,1)$ の候補は，Anokhina ら~\cite{Anokhina12} の $Z$ 因子展開の係数を一つ取り替えたものとして書ける（\S\ref{sec:fig8-21}）．

\begin{table}[t]
\centering
\caption{検証の全体像．○ は spin 表示と genus 表示の両方で正，— は未実施，× は入力が無い．}\label{tab:overview}
\resizebox{\textwidth}{!}{%
\begin{tabular}{llccccc}
\toprule
結び目 & 入力 & $\yng(1)$ & 2 箱 $x=0$ & 2 箱 $x\neq0$ & 3 箱 $x=0$ & 3 箱 $x\neq0$ \\
\midrule
自明な結び目 & -- & 全 $x$ & ○ & ○ & ○ & $|x|\le6$ \\
$T_{2,3},T_{2,5},T_{2,7},T_{3,4}$ & 楕円 Hall 代数 & 全 $x$ & ○ & ○ & ○ & $|x|\le 3$--$6$ \\
$T_{3,5},T_{2,11},T_{3,7},T_{4,5}$ & 楕円 Hall 代数 & 全 $x$ & ○ & $|x|\le10$ & ○ & $|x|\le10$ \\
$K1,K3$ & DAHA & 全 $x$ & ○ & ○ & ○ & $|x|\le2$（$K1$ は $x=\pm6$ も） \\
$K6a,K7,K8,K9$ & DAHA & 全 $x$ & ○ & $|x|\le10$ & ○ & $|x|\le10$ \\
$K10$ & DAHA & 全 $x$ & ○ & $|x|\le10$ & ○ & $|x|\le2$ \\
$\mathbf{4_1}$ & ホモロジー（$\yng(2,1)$ は KN の候補） & 全 $x$ & ○ & ○ & ○ & $|x|\le4$ \\
結び目表 801 個 & ホモロジー & 全 $x$ & × & × & × & × \\
ツイスト $K_p$ & KNTZ の式 & 全 $x$ & 全 $p$ & 注 1 & × & × \\
ダブルツイスト $K_{m,n}$ & KNTZ の式 & 全 $x$ & 注 2 & -- & × & × \\
$\mathbf{10_{139}},\mathbf{10_{152}}$ & quiver ＋ KRSS の予想式 & 全 $x$ & ○ & -- & × & × \\
\bottomrule
\end{tabular}}

\smallskip
\raggedright
注 1：$p$ と $x$ が同符号の向きでは全て正．逆符号は $|p|\le6$，$|x|\le14$ の直接計算で正．\\
注 2：同符号は全て正．異符号は $\min(|m|,n)\le20$ で証明，$|m|\le30$ かつ $n\le40$ で確認．\\
4 箱以上（framing 0）：トーラス結び目 8 個で正（表~\ref{tab:higher}）．$\mathbf{4_1}$ は弱い判定（\S\ref{sec:fig8-four}）．
\end{table}

\subsection{関連する研究}
非 refined の framed knot の LMOV 型の整数性は Mariño--Vafa~\cite{Marino01} に始まり，
Zhu らの一連の研究~\cite{Zhu21,Luo16,Luo21,Wang25,Zhu25} がある．
Zhu~\cite{Zhu21} の ``refined LMOV conjecture'' の refinement は KN の refined（$t\neq q$）とは別の意味であり（同論文の Remark），本論文の主張とは対象が異なる．
refined の open BPS 不変量の非負性は，単一のブレーン・対称表現について Cheng--Sułkowski~\cite{Cheng21}（framing は固定）と
Banerjee--Ishtiaque--Jeong~\cite{Banerjee26} が論じている．
knots--quivers 対応~\cite{Kucharski17,Kucharski17b,Ekholm18} では，framing の変化は quiver 行列のずらし $C\mapsto C+f$ として現れ，
対称表現の refined LMOV 不変量の整数性が quiver の構造から従う．
本論文の調べた範囲では，KN の $\hf_\rho$ を framing をつけて作り，非対称な $\rho$（$\yng(2,1)$ を含む）まで spin 表示と genus 表示の正整数性を調べた研究は見当たらない．
ケーブルの refined 不変量の定義は Cherednik--Danilenko~\cite{Cherednik14,Cherednik15} による．
その明示例は $K1$ の $\yng(1),\yng(2)$ などに限られ，$|\lambda|=3$ の値と refined BPS 不変量は文献に見当たらない．

\subsection{構成}
\S\ref{sec:prelim} で変数と KN の refined LMOV 不変量を復習する．
\S\ref{sec:framing} で framed knot への拡張と，$x$ によらない構造を述べる．
\S\ref{sec:inputs} で入力（トーラス結び目，ケーブル，ホモロジー，予想式）と計算法を説明する．
\S\ref{sec:one}，\S\ref{sec:two}，\S\ref{sec:three}，\S\ref{sec:higher} で，それぞれ $\yng(1)$，2 箱，3 箱，4 箱以上の検証を述べる．
\S\ref{sec:discussion} で未解決の問題を議論する．
付録に計算法の詳細，正規化の導出，非代数的なケーブル，データとコードをまとめる．

主張の確からしさを区別するため，解析的に導いたものを定理・命題・補題とし，計算機で確かめたものは「計算結果」として範囲とともに書く．
有限個の厳密な計算機計算（多項式の係数の符号の判定や線形計画法の証明書の検証）を含む証明は，定理の見出しに「計算機援用」と記す．
予想式を入力とする主張は，その旨を文中に明記する．

\section{準備}\label{sec:prelim}

\subsection{変数}
KN~\cite{Kameyama17} に従い，$q=e^{i\epsilon_1}$，$t=e^{i\epsilon_2}$ を $\Omega$ 背景 $(z_1,z_2)\mapsto(qz_1,t^{-1}z_2)$ の変数とし，
$SU(N)$ の不変量の large $N$ 極限を $a=t^N$ で表す．
太字の変数 $(\ba,\bq,\bt)$ を
\begin{equation}\label{eq:cov}
a=-\ba^2\bt,\qquad q^{1/2}=-\bq\bt,\qquad t^{1/2}=\bq
\end{equation}
で定める（KN の式 (cov)）．これは Dunfield--Gukov--Rasmussen（DGR）~\cite{Dunfield05} の superpolynomial の変数と同じであり，
$t=q$ は $\bt=-1$ に当たる．Aganagic--Shakirov~\cite{Aganagic11} の太字の変数も同じである．
Cherednik--Danilenko~\cite{Cherednik14} の DAHA の変数 $(q,t,a)$ は KN の $(q,t)$ と同じで，$a_{\mathrm{DAHA}}=-a$ である．
半整数冪の分枝は \eqref{eq:cov} で固定する．$\hf_\rho$ の全体の符号は規約であり，最低次の項が正になる向きにとる．

\paragraph{結び目の向き．}
結び目の向き（鏡像のどちらをとるか）は DGR に従う．
トーラス結び目 $T_{m,n}$（$m,n>0$）は，三葉結び目 $T_{2,3}$ の reduced superpolynomial が $\ba^2\bq^{-2}+\ba^2\bq^2\bt^2+\ba^4\bt^3$ になる向きにとる（この向きでは交差は全て負である）．
楕円 Hall 代数では，$T_{m,n}$ は傾き $(m,-n)$ の元に当たる．
結び目表の結び目は KnotInfo~\cite{KnotInfo} のブレイドの向きのまま用い，ツイスト結び目とダブルツイスト結び目は KNTZ~\cite{Kameyama19} の向き（$K_{1,1}=\mathbf{3_1}$ など）にとる．
鏡像 $K^\ast$ の superpolynomial は，$K$ のものの $\ba,\bq,\bt$ を逆数にしたもの（単項式倍を除く）である．
framing $x>0$ は Mariño--Vafa の $p>0$ に当たる（命題~\ref{prop:MV}）．

\subsection{Macdonald 多項式と自明な結び目}
$P_\lambda(u;q,t)$ を monic な Macdonald 多項式とし（$u=(u_1,u_2,\dots)$ は対称関数の変数．framing の $x$ と区別する），内積 $\langle p_\lambda,p_\mu\rangle_{q,t}=\delta_{\lambda\mu}z_\lambda\prod_i(1-q^{\lambda_i})/(1-t^{\lambda_i})$ に関する
$g_\lambda=(q/t)^{|\lambda|}/\langle P_\lambda,P_\lambda\rangle_{q,t}$ を用いる．
自明な結び目の不変量は主特殊化
\[
P_\lambda(1,t,\dots,t^{N-1})=\prod_{s\in\lambda}\frac{t^{l'(s)}-a\,q^{a'(s)}}{1-q^{a(s)}t^{l(s)+1}}\qquad(a=t^N)
\]
を対称化したものである．$a(s),l(s)$ は腕と脚の長さ，$a'(s),l'(s)$ は余腕と余脚の長さである．
$n(\lambda)=\sum_i(i-1)\lambda_i$ とおき，
\[
T_\lambda=q^{n(\lambda')}\,t^{-n(\lambda)}
\]
を refined の framing の固有値と呼ぶ．$t=q$ では $q^{\sum_{s\in\lambda}c(s)}$（$c(s)$ は content）になる．

\subsection{refined CS 不変量}
KN と同じく，unreduced の refined CS 不変量を $\overline{\rCS}_\lambda(K;a,q,t)$，reduced の不変量を $\rCS_\lambda(K;a,q,t)$ と書く．両者は
\[
\overline{\rCS}_\lambda(K;a,q,t)=\overline{\rCS}_\lambda(\bigcirc;a,q,t)\,\rCS_\lambda(K;a,q,t)
\]
で結ばれる．ここで $\bigcirc$ は自明な結び目である．
トーラス結び目では，Aganagic--Shakirov の $S,T$ 行列による定義~\cite{Aganagic11}，Cherednik の DAHA による定義~\cite{Cherednik11}，
Gorsky--Neguţ~\cite{Gorsky13b} の楕円 Hall 代数による表示が知られている．
本論文では後者で計算する（\S\ref{sec:inputs}）．
これらの定義は単項式倍を除いてしか値を決めない．KN の BPS 抽出は単項式倍で結果が変わるので，単項式の正規化を \S\ref{sec:normalization} で固定する．
$\U(N)$ と $\SU(N)$ の違いは framing と linking number の補正だけで，本論文では区別しない（KN の脚注と同じ扱い）．

\subsection{refined reformulated invariants}
KN は，$q$-brane と $\bar t$-brane の二つの配位について，幾何学的遷移の等式
\begin{align}
\sum_\lambda \overline{\rCS}_\lambda(K;a,q,t)\,g_\lambda(q,t)\,P_\lambda(u;q,t)
&=\exp\Bigl(\sum_{d=1}^\infty\sum_\mu \frac1d\,\frac{f^q_\mu(K;a^d,q^d,t^d)}{q^{\frac d2}-q^{-\frac d2}}\,s_\mu(u^d)\Bigr),\label{eq:GT1}\\
\sum_\lambda \overline{\rCS}_\lambda(K;a,q,t)\,P_{\lambda^T}(-u;t,q)
&=\exp\Bigl(\sum_{d=1}^\infty\sum_\mu \frac1d\,\frac{f^{\bar t}_\mu(K;a^d,q^d,t^d)}{t^{-\frac d2}-t^{\frac d2}}\,s_\mu(u^d)\Bigr)\label{eq:GT2}
\end{align}
を提案した．
ここから定まる refined reformulated invariants $f^q_\mu$, $f^{\bar t}_\mu$ は，KN が定めた $t$ だけに依存する行列 $M_{\mu\rho}(t)$（$|\mu|=|\rho|$）により
\begin{equation}\label{eq:f-to-hatf}
f^q_\mu(K;a,q,t)=\sum_\rho M_{\mu\rho}(t)\,\wh f_\rho(K;a,q,t),\qquad f^{\bar t}_\mu(K;a,q,t)=\sum_\rho M_{\mu\rho}(q^{-1})\,\wh f_\rho(K;a,q,t)
\end{equation}
と書け，二つの式が同じ $\wh f_\rho$ を与えることが KN の枠組みの整合性になる．
$f^q_\mu$ の尺度は KN の付録 B の式に合わせる．本文の \eqref{eq:GT1} をそのまま展開した値とは $(q/t)^{|\mu|/2}$ だけ違う（$u$ を $(t/q)^{1/2}u$ に取り替える尺度の違い）．$f^{\bar t}_\mu$ は \eqref{eq:GT2} のとおりで，この違いは無い．

KN の予想は，$a$ 次数を $\pm\frac12$ ずらした後で $\wh f_\rho$ が
\begin{equation}\label{eq:hat-f-2}
\wh f_\rho(K;a,q,t)=\sum_{g,\beta,J_r}(-1)^{2J_r}\,\wh N_{\rho,g,\beta,J_r}(K)\,(q^{\frac12}-q^{-\frac12})^{g}(t^{-\frac12}-t^{\frac12})^{g}\Bigl(\frac qt\Bigr)^{J_r-\frac\beta2}a^{\beta}
\end{equation}
の形をとり，$\wh N_{\rho,g,\beta,J_r}(K)$ が非負整数になる，というものである．
ここで $g$ は M2-brane の曲線の genus，$\beta$ は $\mathbb P^1$ の次数，$J_r$ は M2-brane のモジュライ空間のコホモロジーの右 spin である．
KN の議論では，BPS 状態は $J_r$ のほかに追加の $U(1)_R$ 対称性の電荷 $S_R$ をもち，$\wh N_{\rho,g,\beta,J_r}$ は $S_R$ が BPS 状態に自明に作用するという予想の下での縮退度である．
$a$ 次数のずらしは，$|\rho|=1$ で $a^{-1/2}$，$|\rho|=3$ で $a^{1/2}$ とする．
\eqref{eq:cov} の変数では，$\wh f_\rho$ は
\begin{equation}\label{eq:bold-form}
\wh f_\rho=\sum_{g,\beta,F}\wh{\mathbf N}_{\rho,g,\beta,F}\,(\bq\bt-\bq^{-1}\bt^{-1})^{g}(\bq-\bq^{-1})^{g}\,\ba^{2\beta}\,\bt^{F}
\end{equation}
と展開される（KN の (bold-form)）．
KN の BPS 状態の数え上げでは $\wh{\mathbf N}_{\rho,g,\beta,F}=\sum_{2(J_r-S_R)=F}(-1)^{2S_R}\wh N_{\rho,g,\beta,J_r,S_R}$ であり，
$S_R$ が自明に作用するという予想の下でだけ $\wh{\mathbf N}_{\rho,g,\beta,F}=\wh N_{\rho,g,\beta,J_r}$（$F=2J_r$）となる．
本論文で計算し判定するのは $\wh{\mathbf N}_{\rho,g,\beta,F}$ である．
KN にならい，\eqref{eq:hat-f-2}，\eqref{eq:bold-form} の展開を M2-brane の genus 展開と呼ぶ．以下では $G=(\bq\bt-\bq^{-1}\bt^{-1})(\bq-\bq^{-1})$ と略記する．
genus 展開が整数係数で閉じるのは，$\ba^{2\beta}$ の各行が
\begin{equation}\label{eq:sigma}
\sigma:\ (e_\ba,e_\bq,e_\bt)\mapsto(e_\ba,-e_\bq,e_\bt-e_\bq)
\end{equation}
で不変で，$\bq$ の指数が偶数のときである．
KN は，genus 展開をする前の素朴な変数変換 \eqref{eq:cov} の係数
\[
\wh f_\rho(K;a=-\ba^2\bt,\,q=\bq^2\bt^2,\,t=\bq^2)=\sum_{i,j,k}\wh{\mathbf N}_{\rho;i,j,k}(K)\,\ba^{2i}\bq^{2j}\bt^{k}
\]
も非負整数になることを観察した．
本論文では，この二つを refined LMOV 不変量の二つの表示として区別し，
$\wh{\mathbf N}_{\rho;i,j,k}$ を \emph{spin 表示}，$\wh{\mathbf N}_{\rho,g,\beta,F}$ を \emph{genus 表示} と呼ぶ（どちらも KN の記号）．
名前の根拠は KN の BPS 状態の構成にある．
\begin{itemize}
\item genus 表示：KN は BPS 状態を $\bigoplus C_{\mu\sigma\rho}\,\mathcal{H}_{\sigma,g}\otimes A_{\rho,g,\beta,J_r,S_R}$ と分解し，$\wh N_{\rho,g,\beta,J_r,S_R}=\dim A_{\rho,g,\beta,J_r,S_R}$ と定める．
因子 $G^g$ は M2-brane の曲線 $\Sigma_{g,h}$ の bulk の Jacobian $H^*((T^2)^g)$ の寄与であり，$g$ はその genus である．
Labastida--Mariño--Vafa~\cite{Labastida01} の $\wh N_{R,g,Q}$ の refined 版にあたる．
\item spin 表示：refined index $\mathrm{Tr}\,(-1)^Fq^{J_1-S_R}t^{S_R-J_2}$（$J_1,J_2$ は $\Omega$ 背景の二つの複素平面 $z_1,z_2$ の回転の生成元）の単項式は，\eqref{eq:cov} で $\pm\bq^{2(J_1-J_2)}\bt^{2(J_1-S_R)}$ になる．
したがって $\wh{\mathbf N}_{\rho;i,j,k}$ は，$\beta=i$ と spin の固有値 $J_1-J_2=j$，$2(J_1-S_R)=k$ を固定した BPS 状態の符号付きの数である．
$J_1-J_2$ は非 refined の index $q^{J_1-J_2}$ の冪で，Ooguri--Vafa~\cite{Ooguri99} が $N_{R,Q,s}$ の spin $s$ と呼んだものにあたる．
\end{itemize}
spin 表示は spin の既約表現ではなく固有値（重み）でラベルされた，符号付きの数である（OV の $N_{R,Q,s}$ も同じ）．その非負性は自明ではない．
genus のラベルは，KN の形式的な BPS 状態の分解の中で意味をもつ．
なお knots--quivers 対応の文献~\cite{Ekholm18} では，genus 展開をしない型の不変量を「LMOV invariants」と呼ぶ．本論文では KN にならい，LMOV 不変量は genus 表示を指す．
$f^q$, $f^{\bar t}$ の係数は refined LMOV 不変量ではなく，付録で補足として扱う．

\begin{definition}[正値性]\label{def:positivity}
$\wh f_\rho$ が\emph{正}であるとは，全体の符号を一つ選ぶと，spin 表示 $\wh{\mathbf N}_{\rho;i,j,k}$ と genus 表示 $\wh{\mathbf N}_{\rho,g,\beta,F}$ の全てが非負整数になることをいう．
\end{definition}
KN の予想は，これに加えて $(\rho,g,\beta)$ を固定したときの $2J_r$ の偶奇がそろうことを含む．本論文の正値性はこの偶奇の条件を含まない（偶奇は \S\ref{sec:one} で別に調べる）．

$\yng(1)$ では $\hf_{\yng(1)}=(a^{1/2}-a^{-1/2})\rCS_{\yng(1)}$ である．
したがって，$\rCS_{\yng(1)}$ を太字の変数で書いた $P$ に対し，$\hf_{\yng(1)}$ は（a 次数をずらした後）$(1+\ba^2\bt)\ba^{-2}\bt^{-1}P$ になる．
$P$ が HOMFLY-PT ホモロジーの Poincaré 多項式として解釈されるなら，$\yng(1)$ の正値性は自然に期待される．

\section{framed knot への拡張}\label{sec:framing}

\subsection{定義}
% [F1]
\begin{definition}\label{def:framed}
分割 $\lambda$ に対し，refined の framing の因子を
\[
\theta_\lambda(q,t)=\bigl(-(q/t)^{1/2}\bigr)^{|\lambda|}\,q^{n(\lambda')}\,t^{-n(\lambda)}
=(-1)^{|\lambda|}\prod_{s\in\lambda}q^{a(s)+\frac12}\,t^{-l(s)-\frac12}
\]
とおく．framing $x\in\mathbb{Z}$ の refined CS 不変量を $\overline{\rCS}^{(x)}_\lambda=\theta_\lambda^x\,\overline{\rCS}_\lambda$ で定め（reduced の不変量でも同じ因子 $\rCS^{(x)}_\lambda=\theta_\lambda^x\rCS_\lambda$），
これを \eqref{eq:GT1}，\eqref{eq:GT2} に代入して得る量を $f^{q,(x)}_\mu$，$f^{\bar t,(x)}_\mu$，$\hf^{(x)}_\rho$，$\hbN^{(x)}_{\rho;i,j,k}$，$\hbN^{(x)}_{\rho,g,\beta,F}$ と書く．
$x=0$ が KN の場合である．
\end{definition}

$\theta_\lambda$ は係数 $\pm1$ の単項式であり，plethystic 対数の中の Adams 作用 $\overline{\rCS}_\lambda\mapsto\overline{\rCS}_\lambda(a^k,q^k,t^k)$ では
$\theta_\lambda(q^k,t^k)^x$ が付く．符号 $(-1)^{|\lambda|}$ は $k$ によらない．
この扱いは Mariño--Vafa~\cite{Marino01} の例
$f^{(p)}_{\yng(2)}=q^pW_{\yng(2)}-\tfrac12\bigl(W_{\yng(1)}^2+(-1)^pW_{\yng(1)}(q^2,a^2)\bigr)$ と同じである（Mariño--Vafa の変数 $\lambda$ を $a$ と書く．以下同じ）．

\subsection{framing の因子の由来}
% [F2]
\begin{proposition}\label{prop:MV}
$t=q$ で $\theta_\lambda^x=(-1)^{|\lambda|x}q^{x\kappa_\lambda/2}$ である．
ここで $\kappa_\lambda=|\lambda|+\sum_i(\lambda_i^2-2i\lambda_i)=2\sum_{s\in\lambda}c(s)$ は Mariño--Vafa の $\kappa_R$ である．
したがって $\overline{\rCS}^{(x)}$ の $t=q$ の極限は，Mariño--Vafa~\cite{Marino01} の framed knot の規則 $W^{(p)}_R=(-1)^{\ell p}q^{p\kappa_R/2}W_R$（$\ell=|R|$ は箱の数）に $x=p$ で一致する．
\end{proposition}
\begin{proof}
$t=q$ で $(q/t)^{1/2}=1$，$n(\lambda')-n(\lambda)=\sum_s c(s)=\kappa_\lambda/2$ である．
変数の向き（$q$ と $q^{-1}$）は，自明な結び目の主特殊化と Mariño--Vafa の量子次元を比べて確かめる．
Mariño--Vafa の $\lambda$ を $a$，$q$ を $q$ とすると両者は単項式倍を除いて一致し，$q\mapsto q^{-1}$ では一致しない（$\lambda\neq\lambda^T$ のとき）．
\end{proof}

\begin{remark}
記号 $\theta_\lambda$ は，色付き HOMFLY-PT 多項式の framing の固有値 $\theta_\lambda=q^{\sum_{s}c(s)}a^{|\lambda|/2}$（Rosso--Jones の公式などで使われるもの）の refined 版として選んだ．
定義~\ref{def:framed} の $\theta_\lambda$ は，$a^{|\lambda|/2}$ を含まず，符号 $(-1)^{|\lambda|}$ と箱あたりの $(q/t)^{1/2}$ を含む点で異なる．
$a^{|\lambda|/2}$ と $(q/t)^{1/2}$ は補題~\ref{lem:adams} により正値性に影響せず，符号は命題~\ref{prop:sign} により必要である．
\end{remark}

\begin{remark}[refined topological vertex との関係]\label{rem:Taki}
Taki~\cite{Taki07} の refined topological vertex の framing の因子は
$f_\mu(t,q)=(-1)^{|\mu|}t^{\|\mu^T\|^2/2}q^{-\|\mu\|^2/2}$ である．
$\|\mu\|^2=2n(\mu')+|\mu|$ を使うと $f_\mu(t,q)=\theta_\mu(q,t)^{-1}$ となる．すなわち Taki の因子は定義~\ref{def:framed} の因子の逆数にちょうど一致し，違いは framing の正の向きの取り方だけである．
本論文では，KN の $f^q_\mu$，$f^{\bar t}_\mu$，$\wh f_\rho$ との混同を避けるため，Taki の $f_\mu$ ではなく $\theta_\lambda$ と書く．
Taki は，Iqbal--Kozçaz--Vafa~\cite{Iqbal07} の因子 $(-1)^{|\nu|}t^{n(\nu)}q^{-n(\nu^T)}$ に $(t/q)^{|\nu|/2}$ を掛ける修正が，
K 理論版の Nekrasov 分配関数と一致するために必要であることを示した．
また Aganagic--Shakirov の refined $T$ 行列の固有値 $T_i=c_Tq^{(\lambda_{R_i}+\beta\rho)^2/2}$（$t=q^\beta$）について，
$T_\lambda/T_\emptyset$ は $SU(N)$ の部分を除いて $(a^{1/2}(q/t)^{1/2})^{|\lambda|}\,q^{n(\lambda')}t^{-n(\lambda)}$ に等しい．
箱あたりの単項式を除けば，framing の 1 単位は refined の $T$ 行列の作用である．
\end{remark}

% [F6]
\begin{proposition}\label{prop:bold}
太字の変数では
\[
\theta_\lambda=\bt^{\,|\lambda|+2n(\lambda')}\,\bq^{\,2(n(\lambda')-n(\lambda))}
\]
であり，符号を含まない．例えば $\theta_{\yng(1)}=\bt$，$\theta_{\yng(2)}=\bq^2\bt^4$，$\theta_{\yng(1,1)}=\bq^{-2}\bt^2$，$\theta_{\yng(2,1)}=\bt^5$ である．
\end{proposition}
\begin{proof}
\eqref{eq:cov} から $-(q/t)^{1/2}=\bt$，$q=\bq^2\bt^2$，$t=\bq^2$ である．
\end{proof}

\begin{remark}
命題~\ref{prop:bold} により，framing の変化は（色付き）HOMFLY-PT ホモロジーの次数を符号なしでずらす操作として働く．
非 refined の極限 $\bt=-1$ では $\theta_\lambda^x\to(-1)^{(|\lambda|+2n(\lambda'))x}=(-1)^{|\lambda|x}$ となる．
Mariño--Vafa の符号 $(-1)^{\ell p}$ は，refined の次数のずらし $\bt^{|\lambda|x}$ の非 refined での影と読める．
この符号が KN の形に必要であることは命題~\ref{prop:sign} で示す．
\end{remark}

\subsection{$x$ についての構造}
% [F5]
\begin{lemma}\label{lem:adams}
$c$ を係数 1 の単項式とし，Adams 作用で $c\mapsto c^k$ となるものとする（$c=a^{1/2}$, $(q/t)^{1/2}$ など）．
$\overline{\rCS}_\lambda$ を $c^{|\lambda|}\overline{\rCS}_\lambda$ に置き換えると，$\hf_\rho$ は $c^{|\rho|}\hf_\rho$ になる．
\end{lemma}
\begin{proof}
$c^{|\lambda|}$ は対称関数の変数を $u\mapsto cu$ と取り替えることに当たり，$c$ が Adams 作用と両立するので plethystic 対数と可換である．
\end{proof}
太字の変数で $c=(q/t)^{1/2}=-\bt$ のとき，これは spin 表示と genus 表示の $\bt$ の次数を一様にずらし，全体の符号を変えるだけである．
したがって正値性に影響しない．

% [F7]
\begin{proposition}\label{prop:expsum}
$|\rho|=n$ のとき，$\bt^{-nx}\hf^{(x)}_\rho$ は $x$ の指数和である．根は
$n=2$ で $\{\bq^2\bt^2,\bq^{-2},1,-1\}$，$n=3$ で $\{\bq^6\bt^6,\bq^{-6},\bq^2\bt^2,\bq^{-2},\bt^2,1,-1\}$ の 7 個である．
$x$ の偶奇で分ければ，根は係数 1 の単項式になる．
\end{proposition}
\begin{proof}
$\hf_\rho$ は，$\overline{\rCS}_\lambda$ とその Adams の写しの積を固定した有理関数の係数で一次結合したものである．
framing をつけると，各項には $\prod\theta_{\lambda_i}(q^{k_i},t^{k_i})^x$ が掛かる．
これらを太字の変数で書き出すと，$n=2$ では $\theta_{\yng(2)}=\bq^2\bt^4$，$\theta_{\yng(1,1)}=\bq^{-2}\bt^2$，$\theta_{\yng(1)}^2=\bt^2$，$\theta_{\yng(1)}(q^2,t^2)=-\bt^2$ である．
$\bt^{-2x}$ を掛ければ主張の根が得られる．$n=3$ も同様である．
\end{proof}

指数和の形は，ツイスト結び目 $K_p$ の撚りの数 $p$ についての構造（\S\ref{sec:two}）と同じ型である．
framing は「もう一つの撚りの変数」として振る舞う．
$v=\bq^2\bt$ とおくと $\bq^2\bt^2=v\bt$，$\bq^{-2}=\bt/v$ で，根は $\sigma$ が入れ替える $v^{\pm1}$ の対をなす．

% [F8, F9]
\begin{proposition}\label{prop:x-indep}
任意の $x$ について次が成り立つ．
\begin{enumerate}
\item $\theta_\lambda(1/t,1/q)=\theta_{\lambda^T}(q,t)$．したがって，$\rCS_\lambda$ の転置対称性，$q$-brane と $\bar t$-brane の式 \eqref{eq:GT1}，\eqref{eq:GT2} が同じ $\hf$ を与えること，genus 展開が閉じる条件（$\sigma$ 不変性）は，$x$ によらない．
\item 2 箱について，$\hf^{(x)}_\rho$ が Laurent 多項式であることと $\hf^{(0)}_\rho$ が Laurent 多項式であることは同値である．
\end{enumerate}
\end{proposition}
\begin{proof}
(1) は定義から直接従う．
(2) $|\rho|=2$ の $\hf_\rho$ は，四つの入力 $\rCS_{\yng(2)}$，$\rCS_{\yng(1,1)}$，$\rCS_{\yng(1)}^2$，$\rCS_{\yng(1)}(a^2,q^2,t^2)$ の一次結合を，
共通の分母 $2A^2(Q^4-1)(T^4-1)(Q^2T^2-1)$（$(A,Q,T)=(a,q,t)^{1/2}$）で割ったものである．
分母の各既約因子の零点で分子に残る入力はちょうど二つで，その二つの framing の因子は零点の上で一致する（記号計算で確かめた）．
よって，零点での分子の消え方は $x$ によらない．
\end{proof}

% [F10]
\begin{proposition}\label{prop:sign}
定義~\ref{def:framed} の符号 $(-1)^{|\lambda|x}$ を外し，$T_\lambda^x$ だけを掛けたとする．
このとき $x$ が奇数なら，命題~\ref{prop:x-indep} (2) の分母の零点 $Q=i$，$T=i$（$q=-1$，$t=-1$）で残る二つの入力の因子が一致しない．
したがって，一般の入力で $\hf_{\yng(2)}$，$\hf_{\yng(1,1)}$ は Laurent 多項式でなくなる．
$x$ が偶数なら符号は効かない．また，箱あたりの $(q/t)^{1/2}$ を外しても Laurent 性は保たれる．
\end{proposition}
\begin{proof}
符号を外すと，Adams の写し $\rCS_{\yng(1)}(a^2,q^2,t^2)$ の項の因子だけが $(-1)^x$ 倍変わる．この項が残る零点で一致が壊れる．
\end{proof}

\begin{compresult}
三葉結び目，$x=1$ で符号を外すと，$\hf_{\yng(2)}$，$\hf_{\yng(1,1)}$ の分母に $(1+q)(1+t)$ が残る．
$x=2$ と，符号をつけた $x=1$ では Laurent 多項式である．
\end{compresult}

命題~\ref{prop:sign} は，Mariño--Vafa の「符号は整数性に不可欠である」という観察~\cite{Marino01} の refined 版である．
Mariño--Vafa は，符号を開いた弦の平坦座標の取り替え $\hat u\mapsto\hat u+i\pi p$（$V\mapsto(-1)^pV$）から説明した．
この取り替えは，生成関数 $\sum_RW_R\,\mathrm{Tr}_RV$ を新しい変数で展開し直すことにあたり，Chern--Simons 側の係数 $W_R$ に定数の符号 $(-1)^{|R|p}$ を掛ける．
定数の符号は Adams の写し $W_R(q^d,a^d)$ で $d$ 乗されない（Mariño--Vafa の $f^{(p)}_{\yng(2)}$ の $(-1)^pW_{\yng(1)}(q^2,a^2)$）．
したがってこれは補題~\ref{lem:adams} の型の自明な置き換えではなく，BPS 側では $d$ 重被覆の項に $(-1)^{d|R|}$ が付く．
定義~\ref{def:framed} の符号は，これと同じ扱いである．
一方，$\theta_\lambda$ の残りの部分 $(q/t)^{|\lambda|/2}q^{n(\lambda')}t^{-n(\lambda)}$ は refined の $T$ 行列の固有値に由来する（注意~\ref{rem:Taki}）．$T$ 行列の固有値は符号を含まない．
Aganagic--Shakirov~\cite{Aganagic11} は，三葉結び目の前因子の $\sqrt{q/t}$ も平坦座標の取り替えに関係するはずだと述べている．
refined の場合の平坦座標の取り替えを幾何から導くことは，本論文ではしていない．

\subsection{KN の命題の framed 版}
\begin{lemma}[genus 展開の補題]\label{lem:G}
$\ba^{2\beta}$ の各行が $\sigma$ で不変で $\bq$ の指数が偶数（すなわち genus 展開が閉じる）であり，spin 表示 $\hbN_{\rho;i,j,k}$ が非負なら，genus 表示 $\hbN_{\rho,g,\beta,F}$ も非負整数である．
逆は成り立たない．
\end{lemma}
\begin{proof}
$v=\bq^2\bt$ とおくと，各行は $\bt^l\chi_k$（$\chi_k=v^k+v^{-k}$）の非負結合である．
$\chi_k$ を $w=v+v^{-1}$ の多項式とみなす．
$w_0=\bt+\bt^{-1}$ のまわりの Taylor 係数は $d\chi_k/dw=kU_{k-1}(w/2)$ と母関数
\[
\frac{1}{1-w\xi+\xi^2}=\sum_{g\ge0}\frac{G^g \xi^g}{\bigl((1-\bt \xi)(1-\xi/\bt)\bigr)^{g+1}}
\]
から非負である（$G=w-w_0$）．逆が成り立たない例は非代数的なケーブル $K5c$ にある（付録~\ref{app:nonalg}）．
\end{proof}

% [F12, F13]
\begin{proposition}\label{prop:one}
$P\in\mathbb{Z}_{\ge0}[\ba^{\pm1},\bq^{\pm1},\bt^{\pm1}]$ の $\bq$ の指数が偶数で，$P$ が $\sigma$ \eqref{eq:sigma} で不変とする．
$P$ を $\rCS_{\yng(1)}$ とする $\hf_{\yng(1)}$ は，全ての $x$ について，genus 展開が整数で閉じ，spin 表示も genus 表示も非負である．
\end{proposition}
\begin{proof}
$x=0$ では，$\hf_{\yng(1)}=(1+\ba^2\bt)\ba^{-2}\bt^{-1}P$ も $\sigma$ で不変で非負なので，補題~\ref{lem:G} から従う．
一般の $x$ では $\theta_{\yng(1)}^x=\bt^x$ なので，$\hf^{(x)}_{\yng(1)}=\bt^x\hf_{\yng(1)}$ である．これは genus 表示の $F$ を $x$ だけずらす．
\end{proof}

\begin{corollary}\label{cor:one}
$\rCS_{\yng(1)}$ が結び目の reduced HOMFLY-PT ホモロジーの Poincaré 多項式（DGR の正規化）で与えられるなら，任意の結び目・任意の $x$ で $\yng(1)$ は正である．
\end{corollary}
\begin{proof}
係数は次元なので非負である．$\sigma$ 不変性は Gorsky--Hogancamp--Mellit~\cite{Gorsky21} の対称性の定理である．
$\bq$ の指数が偶数であることは結び目のホモロジーの次数の偶奇から従う．
\end{proof}

% [F15]
\begin{proposition}\label{prop:two-necessary}
太字の変数で非負係数をもち，$\bq$ の指数が全て偶数の入力 $\rCS_{\yng(1)}$，$\rCS_{\yng(2)}$ に対し，$\hf_{\yng(2)}$ が Laurent 多項式なら $\rCS_{\yng(2)}|_{q=1}=\rCS_{\yng(1)}|_{q=1}^2$ である．
特に，入力がホモロジーの Poincaré 多項式なら $\dim H^{\yng(2)}=(\dim H^{\yng(1)})^2$ が必要である．
$\yng(1,1)$ についても $t=1$ で同様である．この条件は $x$ によらない．
\end{proposition}
\begin{proof}
$(A,Q,T)=(a,q,t)^{1/2}$ と書く．2 箱の公式の分子を $Q=1$ に制限すると，$\rho=\yng(2),\yng(1,1)$ のどちらでも
\[
2T(A^2-1)^2(T^4-1)\bigl(\rCS_{\yng(1)}^2-\rCS_{\yng(2)}\bigr)
\]
となり，$\rCS_{\yng(1,1)}$ と Adams の写し $\rCS_{\yng(1)}(a^2,q^2,t^2)$ の項は消える（$T=1$ では $\yng(2)$ と $\yng(1,1)$ の役が入れ替わる）．
分母は $Q=1$ で消えるので，Laurent 性には $Q=1$ で $\rCS_{\yng(2)}=\rCS_{\yng(1)}^2$ が必要である．
$Q=1$ は $\bt=-\bq^{-1}$ にあたり，単項式 $\ba^i\bq^j\bt^k$ は $(-1)^k\ba^i\bq^{j-k}$ になる．
$j$ が偶数なら $(-1)^k=(-1)^{j-k}$ なので，さらに $\bq\mapsto-\bq$ とすると全ての項の係数が非負になり，打ち消しは起きない．
$\ba=\bq=1$ とおけば次元の等式を得る．$x$ によらないことは命題~\ref{prop:x-indep} による．
\end{proof}

\begin{remark}
$\mathbf{9_{42}}$ では $\dim H^{\yng(1)}=9$，$\dim H^{\yng(2)}=401$~\cite{Gorsky13} で $401\neq81$ である．
したがって，$t_c$ 次数の選び方によらず，KN の形に乗らない．
$\mathbf{9_{42}}$ が KN の形に乗らないことは KN~\cite{Kameyama17} が指摘し，thick なホモロジーが指数増大性をもたないこと（Gukov--Stošić~\cite{Gukov11} の付録 B）と結びつけた．
指数増大性 $\dim H^{\yng(2)}=(\dim H^{\yng(1)})^2$ は Gukov--Stošić~\cite{Gukov11}，Fuji--Gukov--Sułkowski~\cite{Fuji12} が論じた性質である．
命題~\ref{prop:two-necessary} は，これが KN の形（Laurent 性）の必要条件であることを示す．
\end{remark}

\subsection{全ての framing での判定法}\label{sec:newton}
命題~\ref{prop:expsum} により，固定した結び目について全ての $x$ の正値性を有限個の判定に帰着できる．
$x$ を偶奇と向き（$x\ge x_0$ または $x\le x_0$）に分け，$x=x_0+2y$ について，根の並び $(r_0,r_1,r_2)$ のニュートン形
\[
\bt^{-nx}\hf^{(x)}_\rho=\sum_{k}c_k\,h_{y-k}(r_0,\dots,r_k)
\]
を作る（$h$ は完全斉次対称式）．係数 $c_k$ が全て同じ符号の非負係数の Laurent 多項式なら，$h$ の係数が非負なので，$y\ge0$ の全てで展開の前が非負である．
$\sigma$ 不変性は命題~\ref{prop:x-indep} から全ての $x$ に従い，補題~\ref{lem:G} で展開の後も非負になる．
起点 $x_0$ と 0 の間は直接判定する．

\subsection{正規化との関係}\label{sec:N2-framing}
% [F17]
\begin{proposition}\label{prop:N2-framing}
\S\ref{sec:normalization} の正規化 (N2$'$) の単項式 $M_\lambda=M_\square^{|\lambda|}T_\lambda^F$ は
\[
M_\lambda=\bigl(M_\square(-(q/t)^{1/2})^{-F}\bigr)^{|\lambda|}\,\theta_\lambda^F
\]
と書ける．すなわち，(N2$'$) のうち表現によって変わる部分は，ちょうど $F$ 単位の framing である．
基準の framing を偶数だけずらすことは，$x$ を同じだけずらし，補題~\ref{lem:adams} の一様なずれを加えることに等しい．
奇数だけずらすと Mariño--Vafa の符号が落ち，命題~\ref{prop:sign} により Laurent 性が壊れる．
\end{proposition}
したがって，全ての $x\in\mathbb{Z}$ についての正値性の主張は，$x=0$ をどこに置くかに依存しない．
Aganagic--Shakirov~\cite{Aganagic11} が三葉結び目の値と DGR の superpolynomial の比を「$-6$ 単位の framing と $\sqrt{q/t}$」と書いていることは，絶対値が $F=mn=6$ に合う（符号は framing の正の向きの取り方による）．

\subsection{物理的な描像と予想}\label{sec:physics}
framed knot の refined large $N$ 双対は，次のように組み立てられる．
\begin{itemize}
\item Chern--Simons 側：$\theta_\lambda$ のうち符号を除いた部分は，refined CS 理論の $T$ 行列の作用である（注意~\ref{rem:Taki}）．
Aganagic--Shakirov の構成では，$U(1)_R$ 対称性は 3 次元多様体が半自由な $U(1)$ 作用をもつ（Seifert 多様体である）ことから存在する．
トーラス結び目の framing の変化は，同じ $S$, $T$ 行列の形式の中で扱われる．
\item 弦側：Mariño--Vafa の議論では，framing の整数は $L_K$ の平坦座標の取り方（無限遠での境界条件）に当たる．符号 $(-1)^{|\lambda|x}$ はこの平坦座標の取り替えから来る（命題~\ref{prop:sign} の後の議論）．
\end{itemize}
次の二点は導出していない．
(i) KN の Schwinger の計算と抽出式 \eqref{eq:GT1}，\eqref{eq:GT2} が，全ての $x$ でそのまま適用できること．
(ii) 弦側で $U(1)_R$ の作用が自明なまま保たれること．
これらを仮定として，次を予想する．

% [F19]
\begin{tcolorbox}
\begin{conjecture}[framed refined LMOV 不変量の正値性]\label{conj:A}
トーラス結び目，および $\rCS_\lambda$ が色付き HOMFLY-PT ホモロジーの Poincaré 多項式で与えられる結び目について，
全ての $x\in\mathbb{Z}$ と全ての $\rho$ で $\hf^{(x)}_\rho$ は正（定義~\ref{def:positivity}）である．
\end{conjecture}
\end{tcolorbox}

\section{入力：KN のデータを超える refined 不変量}\label{sec:inputs}

\subsection{トーラス結び目}
トーラス結び目 $T_{m,n}$ の $\overline{\rCS}_\lambda$ は，Gorsky--Neguţ~\cite{Gorsky13b} による楕円 Hall 代数の生成元の作用で計算する．
楕円 Hall 代数と DAHA の対応は Schiffmann--Vasserot~\cite{Schiffmann09} により，ケーブルの場合の対応は Morton--Samuelson~\cite{Morton14} による．
% [I1]
規約を固定するため，KN~\cite{Kameyama17} の補助ファイルにあるトーラス結び目の値 27 件（$T_{2,3}$, $T_{2,5}$, $T_{2,7}$, $T_{2,9}$ の $|\lambda|\le3$ と $T_{3,4}$ の $\yng(1),\yng(2),\yng(1,1)$）と，
KN の付録 C の BPS の表 13 個を，単項式の正規化（\S\ref{sec:normalization}）の下で全係数で再現した．
これは既知の結果の再現であり，以下の計算の規約の確認として用いる．

\subsection{代数的な iterated torus 結び目}
companion $K$ の $(m,n)$-ケーブルを $C_{m,n}K$ と書く．$m$ は巻き数で，$n$ は companion の Seifert framing で測る．
Cherednik--Danilenko~\cite{Cherednik14,Cherednik15} の DAHA による定義では，$C_{m_2,n_2}T_{m_1,n_1}$ の Newton 対は $r=(m_1,m_2)$，$s=(n_1,n_2-m_1n_1m_2)$ であり，
全ての $r_i,s_i>0$（$n_2>m_1n_1m_2$）のとき代数的である．
% [I2]
我々はこの定義を companion の段で線形に分解し，楕円 Hall 代数の演算子の合成として計算した．
すなわち一番外側の段 $C_{m,n}$ について $\Phi_\lambda=P_\lambda[p_k\mapsto P_{k(m,-s')}]\cdot1=\sum_\mu c_{\lambda\mu}(s')P_\mu$ と展開し（$s'$ は Newton 対 $s$ の最後の成分，$p_k$ は冪和），
companion の値 $\rCS_K(\mu)$ との一次結合をとる．
素朴な Rosso--Jones の $t$ 変形が誤りであることは知られており~\cite{Danilenko14,Morozov18}，ここで使う形は定義から直接従うものである．
\begin{compresult}\label{cr:daha}
この計算法は，Cherednik--Danilenko の明示例 10 件（$K1$ の $\yng(1)$ と $\yng(2)$，$K2,K3,K4,K5a,K5b,K5c,K7,K8$ の $\yng(1)$．非代数的なケーブルの名前は付録~\ref{app:nonalg}）と，単項式倍を除いて全係数で一致する．
また，$t=q$ の極限はブレイドから計算した色付き HOMFLY-PT 多項式と一致する．
\end{compresult}

本論文で扱うケーブルを表~\ref{tab:cables} にまとめる．
\begin{table}[h]
\centering
\caption{代数的な iterated torus 結び目．$C_{m,n}K$ は $K$ の $(m,n)$-ケーブル．$b$ はブレイドの本数，$\delta$ は種数型の量（\S\ref{sec:normalization}）．}\label{tab:cables}
\begin{tabular}{llcc}
\toprule
名前 & 結び目 & $b$ & $\delta$ \\
\midrule
$K1$ & $T_{2,3}$ の $(2,13)$-ケーブル & 4 & 8 \\
$K3$ & $T_{2,3}$ の $(2,15)$-ケーブル & 4 & 9 \\
$K6a$ & $T_{2,5}$ の $(2,21)$-ケーブル & 4 & 14 \\
$K9$ & $T_{2,7}$ の $(2,31)$-ケーブル & 4 & 21 \\
$K8$ & $T_{3,4}$ の $(2,25)$-ケーブル & 6 & 18 \\
$K7$ & $T_{2,3}$ の $(3,19)$-ケーブル & 6 & 21 \\
$K10$ & $K1$ の $(2,53)$-ケーブル & 8 & 42 \\
\bottomrule
\end{tabular}
\end{table}

\subsection{単項式の正規化}\label{sec:normalization}
DAHA と楕円 Hall 代数は値を単項式倍を除いてしか決めないが，KN の BPS 抽出は単項式倍で結果が変わる．
そこで次の規則で正規化する．
\begin{itemize}
\item[(N1)] $t=q$ の極限が，framing 0 の reduced 色付き HOMFLY-PT 多項式（DGR の規約）に等しい．
\item[(N2)] $\yng(1)$ の太字の変数での最低の $\bt$ 次数が 0 である．
\item[(N2$'$)] 表現 $\lambda$ の単項式は $M_\lambda=M_\square^{|\lambda|}\,T_\lambda^{F}$ である．$F$ は一番外側のケーブル $C_{m,n}$ では $mn$，トーラス結び目 $T_{m,n}$ では $mn$ である．
\end{itemize}
トーラス結び目では，(N1)(N2) の下で KN の値 27 件が全係数で再現される．
(N2) が決めるのは，箱あたりの $(q/t)^{c}$ の自由度だけである．補題~\ref{lem:adams} により，これは正値性に影響しない．
色付きの場合の単項式 $\pm a^\alpha q^\beta t^\gamma$ のうち，$t=q$ で決まるのは符号，$\alpha$，$\beta+\gamma$ である．
色付きの場合の (N1) は検査として確かめた．
$t=q$ で決まらないのは，表現ごとの $(q/t)^{c_\lambda}$ の自由度だけである．

% [I4]
\begin{hypothesis}\label{hyp:N2}
ケーブルについて，$t=q$ で決まらない表現ごとの単項式 $(q/t)^{c_\lambda}$ は (N2$'$) の形をとる．
すなわち，$|\lambda|$ に比例する部分（正値性に無関係）と，refined の framing の固有値 $T_\lambda^F$ の分だけである．
\end{hypothesis}
命題~\ref{prop:N2-framing} により，仮説~\ref{hyp:N2} は「表現によって変わる部分は Seifert framing の $F$ 単位の framing である」と言い換えられる．
根拠は次の三つである．
(i) 後述の refined 衛星公式（仮説~\ref{hyp:sat}）から従う．
(ii) 非代数的なケーブル $K5c$ では，$\yng(2)$ を $(q/t)^u$ 倍する 1 パラメータ族のうち，Laurent 性と転置対称性を満たすのは (N2$'$) の $u$ だけである．
(iii) 9 個のケーブルの全ての表現で，(N2$'$) の形が計算した値と一致する．

% [I5]
\begin{compresult}\label{cr:M}
(N1)(N2) で決めた $\yng(1)$ の単項式は，計算した全 22 結び目（トーラス結び目 9 個，ケーブル 13 個）で
\[
M_\square=(-q/t)^{b-1}\bigl(a(q/t)^{1/2}\bigr)^{\delta}
\]
に一致する．ここで $b$ はブレイドの本数，$\delta$ はケーブル $C_{m,n}$ の段ごとに $\delta\mapsto m\delta+(m-1)(n-1)/2$ で決まる（$T_{m,n}$ では $\delta=(m-1)(n-1)/2$）．
\end{compresult}

% [I6]
\begin{hypothesis}[refined 衛星公式]\label{hyp:sat}
$\rCS_{C_{m,n}K}(\lambda)=\sum_\mu C_{\lambda\mu}\,\rCS_K(\mu)$ が成り立つ．
係数は，トーラス結び目 $T_{m,n}$（自明な結び目の $(m,n)$-ケーブル）の展開係数を (N2$'$) で正規化した $C_{\lambda\mu}=M_{T_{m,n}}^{|\lambda|}T_\lambda^{mn}c_{\lambda\mu}(n)$ である．
\end{hypothesis}
$t=q$ では仮説~\ref{hyp:sat} は古典的な衛星公式に帰着し，(N1) と両立する．

導出には，計算で確かめた次の二つの式を用いる．
\begin{compresult}\label{cr:tau}
framing 演算子 $\tau:P_\mu\mapsto T_\mu P_\mu$ について，楕円 Hall 代数の傾き $(m,n)$ の元と整数 $j$ に対し
$\tau^jP_\lambda[p_k\mapsto P_{k(m,n)}]\cdot1=P_\lambda[p_k\mapsto P_{k(m,n+jm)}]\cdot1$ が成り立つ（5 つの傾きと $j=1,-1,2$ で，有理数の点で厳密に一致）．
\end{compresult}
\begin{compresult}\label{cr:mellit}
Mellit~\cite[Thm.~3.1]{Mellit17} の Dyck 路の表示によるトーラス結び目の HOMFLY-PT ホモロジーの Poincaré 多項式 $\mathcal{P}_{m,n}$ は，
\[
\mathcal{P}_{m,n}=-\bigl(a(qt)^{-1/2}\bigr)^{\delta}\frac{1}{1-t}\,\frac{(Q_{m,n}\cdot1)[a/q-1]}{a/q-1}
\]
に一致する（$Q_{m,n}$ は Bergeron--Garsia--Leven--Xin~\cite{Bergeron14} の演算子．$(m,n)=(2,3),(3,2),(2,5),(3,4),(4,3),(3,5),(5,3)$ で厳密に一致）．
\end{compresult}

\begin{proposition}\label{thm:M}
仮説~\ref{hyp:sat} と計算結果~\ref{cr:tau}，\ref{cr:mellit} を仮定すると，計算結果~\ref{cr:M} の閉じた式と (N2$'$) の形が成り立つ．
\end{proposition}
\begin{proof}[証明の概略]
ケーブルの部分では，計算結果~\ref{cr:tau}（framing 演算子が傾きを $(m,n)\mapsto(m,n+jm)$ に送ること）を使う．
これと仮説~\ref{hyp:sat} から，$M_{C_{m,n}K}=M_K^m\,M_{T_{m,n}}$ と $F=mn$ が従う．
トーラス結び目の部分は四つの事実からなる．
(I) Gorsky--Neguţ の正規化前の値の最低 $a$ 係数は，有理 $q,t$-Catalan 数で与えられる．これは Gorsky--Neguţ の予想であり，Mellit~\cite{Mellit16} の定理の系として成り立つ．
(II) 最低 $a$ 次数は，Franks--Williams の定理による．
(III) $t=q$ での符号と $q$ の冪は，$q\leftrightarrow q^{-1}$ 対称性から決まる．
(IV) $(q/t)$ の指数は，Mellit の Dyck 路の表示（計算結果~\ref{cr:mellit}）から決まる．
詳細は付録~\ref{app:normalization} に述べる．
\end{proof}

\begin{remark}[Mellit の系の正規化]
Mellit~\cite[Cor.~3.4]{Mellit17} は $\mathcal{P}_{m,n}=(-1)^n(a(qt)^{-1/2})^\delta\frac{1}{1-t}\frac{P_{m,n}(1)[a-1]}{a-1}$ と書き，$P_{m,n}$ を「対応する楕円 Hall 代数の生成元」とする．
同論文の Thm.~3.3 から，この式は $P_{m,n}(1)=-d_-B_{c_{m,n}}d_+(1)$（Carlsson--Mellit の $d_\pm$ の表現の元）と読めば成り立つ．
一方，標準の正規化（Bergeron--Garsia--Leven--Xin の $Q_{m,n}$，Gorsky--Neguţ の $\widetilde P$）の元と読むと，代入 $[a-1]$ では成り立たない．
標準の正規化の元と変数の取り方を変えて調べると，成り立つのは計算結果~\ref{cr:mellit} の形だけである．
命題~\ref{thm:M} の (IV) が決めるのは $M_\square$ の $(q/t)$ の指数だけなので，この点は正値性の結果に影響しない（補題~\ref{lem:adams}）．
\end{remark}

\begin{remark}[Cherednik--Danilenko の tilde 正規化との関係]\label{rem:tilde}
Cherednik--Danilenko~\cite{Cherednik14} は DAHA superpolynomial を，最低の $a$ の冪と $q_{st}^\bullet t_{st}^\bullet$ で割って，$a_{st}=0$ の部分の定数項が 1 になるように正規化する（tilde 正規化）．
標準の変数 $(q_{st},t_{st},a_{st})$ は太字の変数と同じである．
係数が非負の計算値 19 個（トーラス結び目 9 個，ケーブル 10 個）で，本論文の $\yng(1)$ の太字の $P$ は
\[
P=\ba^{2\delta}\bq^{-2\delta}\times(\text{tilde 正規化})
\]
である（計算）．Cherednik--Danilenko の明示例との比も，正のもの（$K1,K3,K5b,K7,K8$）で全て $\ba^{2\delta}\bq^{-2\delta}$ である．
これらの $P$ は全て $\sigma$ で不変で，$\bq$ の指数は偶数である．
\end{remark}

\subsection{ホモロジーを入力とする結び目}
11 交点以下の全素結び目 801 個について，reduced HOMFLY-PT ホモロジーの Poincaré 多項式（Nakagane--Sano~\cite{Nakagane21} のデータ）を $\yng(1)$ の入力とする．
DGR の規約への変換は，$\sigma$ 不変性，$\bt=-1$ での HOMFLY-PT 多項式，トーラス結び目の $\rCS_{\yng(1)}$ との一致で確かめた．

\subsection{2 箱の予想式}
双曲結び目の refined な $\yng(2)$, $\yng(1,1)$ は定義どおりには得られない．そこで文献の予想式を入力とする．
\begin{itemize}
\item ダブルツイスト結び目 $K_{m,n}$ とツイスト結び目 $K_p=K_{p,1}$：Kameyama--Nawata--Tao--Zhang（KNTZ）~\cite{Kameyama19} の，長方形のヤング図で色付けた HOMFLY-PT ホモロジーの Poincaré 多項式の式（refined 版，予想）を用いる．
KNTZ の $K_{m,n}$ は，GNSSS~\cite{Gukov15} の $(2m-1,1,2n-1)$-pretzel 結び目 $K_{s,t}$（$s=m$，$t=n$）と同じ結び目である．
KNTZ の式は GNSSS の四重次数の予想式と一致する（$\yng(1)$，$\yng(2)$，$\yng(1,1)$ について．下の計算結果~\ref{cr:kntz}）．
$t=q$ では KNTZ の式は Fang--Zhou~\cite{Fang26} により証明されている．
KN の $t_c$ 次数への変換は，Gorsky--Gukov--Stošić~\cite{Gorsky13} の次数の関係 $\mathbf{Q}\mapsto\bq^R$，$\mathbf{t}_r\mapsto\bq^{-1}$，$\mathbf{t}_c\mapsto\bt\bq$（$R$ は長方形の行数）による．
変換は，三葉結び目 $K_{1,1}=T_{2,3}$ の $\yng(1)$，$\yng(2)$，$\yng(1,1)$，$\yng(2,2)$ が楕円 Hall 代数の値と全係数一致することで較正した（Nawata ら~\cite{Nawata12} も参照）．
\item thick な結び目 $\mathbf{10_{139}}$，$\mathbf{10_{152}}$：Stošić~\cite{Stosic24} の quiver 型の式と，Kucharski--Reineke--Stošić--Sułkowski（KRSS）~\cite{Kucharski17b} の予想式で $t_c$ 次数を作る．
この方法は，$\mathbf{8_{19}}=T_{3,4}$ で楕円 Hall 代数の値と $\yng(2)$ から $\yng(1,1,1)$ まで全項一致する．
\end{itemize}
\begin{compresult}\label{cr:kntz}
KNTZ の refined の式と GNSSS の四重次数の式は，$K_{m,n}$ の $\yng(1)$，$\yng(2)$ と，ミラー対称性 $\mathbf{t}_r\leftrightarrow\mathbf{t}_c$ を介した $\yng(1,1)$ で，
符号の四つの象限それぞれの $|m|,|n|\in\{1,2,3\}$ で全係数一致する（144 件）．
両方の式は $m,n$ のそれぞれについて（象限ごとに）根が 3 個以下の指数和なので，全ての $(m,n)$ で一致する．
また KNTZ の式の実装は，KNTZ~\cite{Kameyama19} の補助ファイルの値と全係数一致する．
\end{compresult}
これらを入力とする主張は「予想式を入力とした場合の正値性」であり，定義どおりの入力による 3 箱の結果（\S\ref{sec:three}）とは確からしさが異なる．

\subsection{厳密な計算}
$\overline{\rCS}_\lambda$ と $\hf_\rho$ は，有理数の点ではなく，31 ビットの素数を法とする評価と補間で求める．
係数が大きいときは 2〜3 個の素数と中国剰余定理を用いる．
得られた多項式は，補間に使っていない素数の点で検証する．
最大の例 $K10$ の 3 箱は，spin 表示で約 183 万項，係数の最大は約 $1.1\times10^{19}$ である．
計算時間とメモリは付録~\ref{app:computation} にまとめる．

\section{検証 I：基本表現}\label{sec:one}

系~\ref{cor:one} により，ホモロジーを入力とする限り，任意の結び目・任意の framing で $\yng(1)$ は正である．
この節では，それを結び目表で確かめ，$2J_r$ の偶奇の振る舞いを調べる．

% [R1]
\begin{compresult}\label{cr:table}
11 交点以下の全素結び目 801 個（reduced HOMFLY-PT ホモロジーが一つの $\delta$ 次数に集中するものを thin，そうでないものを thick と呼ぶ．thin で交代的なもの 562 個，thin で非交代的なもの 143 個，thick なもの 96 個）で，
$\hf_{\yng(1)}$ の genus 展開は整数で閉じ，spin 表示も genus 表示も非負である．
\end{compresult}

ケーブルの $\yng(1)$ については，注意~\ref{rem:tilde} のとおり，計算した値が非負で $\sigma$ 不変である．したがって命題~\ref{prop:one} を直接適用できる．
一般の代数的な iterated torus 結び目については，Cherednik--Danilenko の予想（\S\ref{sec:conjB}）の下で同じ結論が従う．

\subsection{$2J_r$ の偶奇と parity knot}
KN は，トーラス結び目では $(\rho,g,\beta)$ を固定すると $2J_r$ の偶奇がそろうことを観察し，これを $S_R$ の自明な作用の証拠とした．
$\yng(1)$ ではこの条件を入力の言葉で書ける．

% [F14]
\begin{proposition}\label{prop:parity}
命題~\ref{prop:one} の仮定の下で，genus 表示 $\hbN_{\yng(1),g,\beta,F}$ の $F$（予想の下では $2J_r$）の偶奇が $(g,\beta)$ ごとにそろうことは，次と同値である．
$P$ の $\ba$ の次数が連続する範囲ごとに，$e_\bt+e_\ba/2$ の偶奇が一定である．
$\ba$ の次数の台に隙間が無いとき，これは Chandler--Gorsky~\cite{Chandler22} の parity knot の条件（HOMFLY-PT ホモロジー全体で，DGR の $t$ 次数と $a$ 次数の半分の差の偶奇が一定）と同じである．
\end{proposition}
\begin{proof}
$\hf_{\yng(1)}$ の $\ba^{2\beta}$ の行には $P$ の隣り合う二つの行だけが入る．
$\chi_k$（補題~\ref{lem:G}）の $G^g$ の係数に現れる $\bt$ の指数は $k-g$ と同じ偶奇をもつので，$2J_r\equiv e_\bt-g$ である．
$\ba$ の次数は偶数なので，$\pm e_\ba/2$ の偶奇は同じである．
\end{proof}

\begin{compresult}
801 個の全てで $\ba$ の次数の台に隙間は無く，$2J_r$ の偶奇がそろう結び目は 9 個である．
$\mathbf{3_1},\mathbf{5_1},\mathbf{7_1},\mathbf{9_1},\mathbf{11a_{367}},\mathbf{8_{19}},\mathbf{10_{124}}$（トーラス結び目）と，$\mathbf{10_{125}}$（thick），$\mathbf{10_{140}}$（thin）である．
これらは Chandler--Gorsky の定義による parity knot と一致する．
\end{compresult}

したがって，$2J_r$ の偶奇の一致は $S_R$ の作用の自明さそのものではなく，ホモロジーの parity の性質を反映している．
非トーラス結び目にも偶奇がそろう例がある．
parity knot の条件そのものは Chandler--Gorsky による．BPS 不変量の偶奇との同値性と，11 交点以下の一覧は本論文で与える．

\section{検証 II：2 箱}\label{sec:two}

この節の双曲結び目の入力は予想式（\S\ref{sec:inputs}）であり，主張は全て「その予想式を入力とした場合」のものである．
KN~\cite{Kameyama17} は，GNSSS の式を入力として，ツイスト結び目とダブルツイスト（pretzel）結び目の族のいくつかの結び目で 2 箱を計算し，形 \eqref{eq:hat-f-2} に乗ることを観察した．
本節の定理~\ref{thm:twist}，\ref{thm:pretzel} の新しい点は，族全体（全ての撚りの数）での正値性の証明にある．
トーラス結び目と代数的なケーブルの入力は定義どおりの refined 不変量である．

\subsection{ツイスト結び目}
$K_p=K_{p,1}$ を KNTZ~\cite{Kameyama19} の規約のツイスト結び目とする（$K_1=\mathbf{3_1}$，$K_2=\mathbf{5_2}$，$K_{-1}=\mathbf{4_1}$，$K_{-2}=\mathbf{6_1}$．GNSSS~\cite{Gukov15} の規約と同じ）．
入力は KNTZ の refined の式である（\S\ref{sec:inputs}）．
$X=\ba^2\bt^2$，$Y=\ba^4\bq^4\bt^8$，$\bar Y=\ba^4\bq^{-4}\bt^4$ とおく．

\begin{lemma}\label{lem:twist-expsum}
$\hf_\rho(K_p)$（$|\rho|=2$）は $p$ について根 $\{1,X,X^2,Y,\bar Y\}$ の指数和である（$p<0$ では根を逆数にし，全体に $\bt^2$ が掛かる）．
\end{lemma}
\begin{proof}
KNTZ の式では，$p$ は捻りの因子 $\mathbf{Tw}^{(p)}_\lambda$（$\lambda\subseteq\yng(2)$ または $\lambda\subseteq\yng(1,1)$）に固有値の $p$ 乗としてだけ入る．
太字の変数では，これは $\bigl((\ba\bt)^{2j}(\bt\bq)^{2j(j-1)}\bigr)^p$（$j=0,1,2$）の形である（GNSSS の式でも同じ形）．
入力の根は $P_{\yng(1)}$ で $\{1,X\}$，$P_{\yng(2)}$ で $\{1,X,Y\}$，$P_{\yng(1,1)}$ で $\{1,X,\bar Y\}$ である．
$\hf$ は入力の積と Adams 作用の一次結合を，命題~\ref{prop:two-necessary} の分母で割ったものである．
\end{proof}

% [R2]
\begin{theorem}[計算機援用]\label{thm:twist}
KNTZ の refined の式を入力とすると，全ての $p\neq0$ で $\hf_{\yng(2)}(K_p)$，$\hf_{\yng(1,1)}(K_p)$ は正である．
\end{theorem}
\begin{proof}[証明の概略]
並び $(1,X,X^2,Y,\bar Y)$ のニュートン形 $s_{1+n}=\sum_{k=0}^{4}c_k\,h_{n-k}(m_0,\dots,m_k)$ を作る．
$c_k$ は $s_1,\dots,s_5$（$|p|=1,\dots,5$ の値）の整数係数の一次結合なので，Laurent 多項式である．
両方の $\rho$，両方の $p$ の符号で，$c_0,\dots,c_4$ の全係数が同じ符号であることを有限個の厳密な計算で確かめる．
$h$ の係数は非負なので，全ての $p$ で展開の前が一様な符号になる．
$\sigma$ は根の集合を保つ（$Y\leftrightarrow\bar Y$）ので，$\sigma$ 不変性は $|p|\le5$ から全ての $p$ に従い，補題~\ref{lem:G} で展開の後も非負になる．
係数 $c_k$ と判定のデータは付録~\ref{app:data} に置く．

別証明として，$-\hf_\rho(K_{\pm n})$ を根の部分集合ごとの完全斉次対称式の非負結合で表す証明書もある（線形計画法で求め，有理数で厳密に検証した）．
こちらは $c_2$ の符号を使わない．$c_2$ が一様な符号をもつ構造的な理由は分かっていない．
\end{proof}

\begin{compresult}
$|p|\le30$ の 120 件で，spin 表示と genus 表示の両方で負の項は 0 である．
閉じた形は，作るのに使っていない $|p|=6,\dots,12,15,20,30$ で直接計算と全係数一致する．
\end{compresult}

\begin{remark}[単峰性]
$v=\bq^2\bt$ とおくと，$c_3$ と $c_4$ は構造から正であり，$\hf(K_p)$ は全ての $p$ で $v$ についての $\mathfrak{sl}_2$ 指標の正係数の和（単峰）になる（$c_0,\dots,c_2$ は有限個の判定）．
$\yng(1)$ の単峰性は Gorsky--Hogancamp--Mellit~\cite{Gorsky21} の hard Lefschetz の性質から全ての結び目で従う．
2 箱の単峰性の構造的な理由は未解明である．
\end{remark}

\subsection{ダブルツイスト結び目}
% [R4]
\begin{theorem}[計算機援用]\label{thm:pretzel}
KNTZ の refined の式を入力とすると，ダブルツイスト結び目 $K_{m,n}$ の 2 箱は次の範囲で正である．
\begin{enumerate}
\item $m,n$ が同符号の全ての $(m,n)$．
\item $m<0<n$ で $\min(|m|,n)\le20$ の全ての $(m,n)$．
\end{enumerate}
\end{theorem}
\begin{proof}[証明の概略]
$\epsilon=(\mathrm{sign}\,m+\mathrm{sign}\,n)/2$ として $\hf(K_{m,n})=\bt^{2(1-\epsilon)}F(m,n)$ と書くと，$F$ は $\mathbb{Z}^2$ 全体で一つの指数和である．
同符号の領域では，2 変数のニュートン形の係数が全て同じ符号になる．
異符号の領域では 2 変数のニュートン形は一様にならない．そこで直線ごとの 1 変数のニュートン形で $\min(|m|,n)\le20$ を示す．
\end{proof}

\begin{compresult}
$|m|,|n|\le20$ の 820 点と，$21\le|m|\le30$・$21\le n\le40$ の 200 点で，両方の $\rho$ についてspin 表示と genus 表示の両方で負の項は 0 である．
\end{compresult}

異符号で $|m|,n>20$ の領域は，くさび形への分割と鏡像型の恒等式で調べる．
くさびの方向のニュートン係数 $d_2,d_3,d_4$ の非正値性は，線形計画法の証明書で示した．
残りは，1 変数の三つの数列の非正値性に帰着する．三つとは，対角線 $F(-u,u)$ と，両方の $\rho$ の $d_1(u)$ である．
これらは $u\le25$ で成り立つが，一般には未証明である．

\subsection{thick な非トーラス結び目}
% [R5]
\begin{compresult}\label{cr:thick}
Stošić の quiver と KRSS の予想式から作った $\mathbf{10_{139}}$，$\mathbf{10_{152}}$ の $\yng(2)$，$\yng(1,1)$ の入力で，2 箱は正である．
次元は $441=21^2$，$841=29^2$ で，命題~\ref{prop:two-necessary} の必要条件を満たす．
$t=q$ の極限は，$(2,s)$-ケーブルから逆に解いた $\yng(2)$，$\yng(1,1)$ の色付き HOMFLY-PT 多項式と一致する．
\end{compresult}
これらは，命題~\ref{prop:two-necessary} の必要条件を満たす thick な非トーラス結び目で 2 箱の正値性が成り立つ例である．
knots--quivers 対応の DT 不変量の正値性からは従わない．
対称表現だけの quiver 級数からは，$\yng(1,1)$ の入力を要する $\hf_{\yng(2)}$，$\hf_{\yng(1,1)}$ は出てこない．

\begin{remark}[正値性は自動的ではない]
$\mathbf{5_2}$ で $t_c$ 次数の割り当てを別の 12 通りに変えると，いずれも負の項が出る．
正値性は入力の細かい次数に依存する，非自明な性質である．
\end{remark}

\subsection{全ての framing}
% [R7]
\begin{theorem}[計算機援用]\label{thm:two-all-x}
自明な結び目，$T_{2,3}$，$T_{2,5}$，$T_{2,7}$，$T_{3,4}$，$\mathbf{4_1}$，$K1$，$K3$ の 2 箱は，全ての $x\in\mathbb{Z}$ で正である．
$\mathbf{4_1}$ の入力は KN のデータである．
\end{theorem}
\begin{proof}
\S\ref{sec:newton} の方法による．各結び目・各 $\rho$・$x$ の偶奇と向きの四つの組で，ニュートン形の係数が一様な符号になる起点 $x_0$ がある．
起点と 0 の間の $x$ は直接判定した．
\end{proof}
負の向きの起点は，結び目が複雑になるほど 0 から遠い（$T_{2,2k+1}$ で $-(2k+2)$，$K1$ で $-20$，$K3$ で $-22$）．この意味は未解明である．

% [R8]
\begin{theorem}[計算機援用]\label{thm:twist-x}
KNTZ の refined の式を入力とすると，ツイスト結び目の 2 箱は，$p$ と $x$ が同符号（$p>0,x\ge0$ または $p<0,x\le0$）の全ての $(p,x)$ で正である．
\end{theorem}
\begin{proof}[証明の概略]
$\bt^{-2x}\hf^{(x)}_\rho(K_p)$ は $p$ と $x$ の両方について指数和である（補題~\ref{lem:twist-expsum}，命題~\ref{prop:expsum}）．
テンソル積のニュートン形の係数を $|p|=1,\dots,5$ と 3 点の $x$ から作り，同符号の向きで一様な符号であることを確かめる．
\end{proof}

\begin{compresult}
逆符号の領域では，$|p|\le6$，$|x|\le14$ の 696 件で負の項は 0 である．この領域の証明書は得られていない．
\end{compresult}

\begin{compresult}[framing を振った直接の判定]\label{cr:framing-batch}
次の 1821 件（$(x,\rho)$ の組）で，spin 表示と genus 表示の両方で負の項は 0 である（9 件は $\hf=0$）．
3 箱までは，自明な結び目と $T_{2,3}$（$|x|\le6$），$\mathbf{4_1}$ と $T_{2,5}$（$|x|\le4$），$T_{2,7}$ と $T_{3,4}$（$|x|\le3$），
$T_{3,5},T_{2,11},T_{3,7},T_{4,5}$ と $K6a,K7,K8,K9$（$|x|\le10$），$K3$（$|x|\le2$），$K1$（$|x|\le2$ と $x=\pm6$）である．
$K10$ は 2 箱まで $|x|\le10$，3 箱は $|x|\le2$ である．
ツイスト結び目 $K_{\pm1},\dots,K_{\pm6}$（$|x|\le4$）は 2 箱までである．
最大の例は $K10$ の 3 箱，$x=2$ で，spin 表示で約 189 万項である．
\end{compresult}

\begin{table}[h]
\centering
\caption{framing を振った判定．$|\rho|=1,2,3$ の列は調べた framing $x$ の範囲，件数は判定した $(x,\rho)$ の組の数（「各」は 1 結び目あたり．$K10$ の 3 箱の $x=0$ は表~\ref{tab:three} の判定で，件数に含めない），負（spin/genus）は spin 表示と genus 表示の負の項の数の合計（全体の符号は良い方）．自明な結び目の 9 件は $\wh f=0$．}\label{tab:framing}
\resizebox{\textwidth}{!}{% scripts/framed_knot/paper_framing_table.py が logs/framed_knot/framing_batch.log と data/framed_knot/framed/ から生成。
% 判定件数の合計 1821（重なりは 1 件と数える）、\hat f = 0 は 9 件。
\begin{tabular}{lccccc}
\hline
結び目 & $|\rho|=1$ & $|\rho|=2$ & $|\rho|=3$ & 件数 & 負（spin/genus） \\
\hline
自明な結び目 & $|x|\le6$ & $|x|\le6$ & $|x|\le6$ & 78 & 0/0 \\
$T_{2,3}$ & $|x|\le6$ & $|x|\le6$ & $|x|\le6$ & 78 & 0/0 \\
$T_{2,5}$, $\mathbf{4_1}$ & $|x|\le4$ & $|x|\le4$ & $|x|\le4$ & 各 54 & 0/0 \\
$T_{2,7}$, $T_{3,4}$ & $|x|\le3$ & $|x|\le3$ & $|x|\le3$ & 各 42 & 0/0 \\
$T_{3,5}$, $T_{2,11}$, $T_{3,7}$, $T_{4,5}$ & $|x|\le10$ & $|x|\le10$ & $|x|\le10$ & 各 126 & 0/0 \\
$K1$ & $|x|\le2$ & $|x|\le2$ & $|x|\le2$, $x=\pm6$ & 36 & 0/0 \\
$K3$ & $|x|\le2$ & $|x|\le2$ & $|x|\le2$ & 30 & 0/0 \\
$K6a$, $K7$, $K8$, $K9$ & $|x|\le10$ & $|x|\le10$ & $|x|\le10$ & 各 126 & 0/0 \\
$K10$ & $|x|\le10$ & $|x|\le10$ & $|x|\le2$ & 75 & 0/0 \\
$K_p$（$1\le|p|\le6$） & $|x|\le4$ & $|x|\le4$ & -- & 各 27 & 0/0 \\
\hline
\end{tabular}
}
\end{table}

\section{検証 III：3 箱}\label{sec:three}

3 箱には非長方形の表現 $\yng(2,1)$ が要る．
定義どおりの $\rCS_{\yng(2,1)}$ が得られるのは，トーラス結び目と iterated torus 結び目（DAHA，楕円 Hall 代数）である．
この節の主な結果は，これらの結び目についてのものである．

\subsection{トーラス結び目と代数的な iterated torus 結び目}
% [R10, R11]
\begin{compresult}\label{cr:three}
表~\ref{tab:three} の 15 個の結び目で，正規化 (N1)(N2$'$)（\S\ref{sec:normalization}）の下で，$|\rho|\le3$ の全 6 つの $\rho$ について $\hf_\rho$ は正（framing 0）である．
内訳はトーラス結び目 7 個，代数的な iterated torus 結び目 7 個，$\mathbf{4_1}$ である．
同じ全体の符号で，spin 表示 $\hbN_{\rho;i,j,k}$ と genus 表示 $\hbN_{\rho,g,\beta,F}$ がともに非負整数になる．
\end{compresult}

\begin{table}[h]
\centering
\caption[3 箱までの判定（framing 0）]{3 箱までの判定（framing 0）．$\checkmark(n)$ は spin 表示と genus 表示の両方で負の項が 0 で，$n$ は spin 表示の項数．
$T_{2,3}$，$T_{2,5}$ と $\mathbf{4_1}$ は KN の表の再現である．$\mathbf{4_1}$ の $\yng(2,1)$ は KN の候補を用いた．}\label{tab:three}
\resizebox{\textwidth}{!}{% scripts/framed_knot/paper_summary.py が生成。\checkmark(n) は spin 表示と genus 表示の両方で負の項 0（n は spin 表示の項数）。それ以外は spin 表示の負/項数; genus 表示の負/項数（前後それぞれで全体の符号は良い方）。framing 0。
\begin{tabular}{llcccccc}
\hline
結び目 & 分類 & $\yng(1)$ & $\yng(2)$ & $\yng(1,1)$ & $\yng(3)$ & $\yng(2,1)$ & $\yng(1,1,1)$ \\
\hline
$T_{2,3}$ & torus & \checkmark\,(6) & \checkmark\,(26) & \checkmark\,(41) & \checkmark\,(128) & \checkmark\,(236) & \checkmark\,(282) \\
$T_{2,5}$ & torus & \checkmark\,(10) & \checkmark\,(112) & \checkmark\,(144) & \checkmark\,(753) & \checkmark\,(986) & \checkmark\,(1083) \\
$T_{3,4}$ & torus & \checkmark\,(21) & \checkmark\,(429) & \checkmark\,(504) & \checkmark\,(3249) & \checkmark\,(3802) & \checkmark\,(4053) \\
$T_{3,5}$ & torus & \checkmark\,(30) & \checkmark\,(766) & \checkmark\,(867) & \checkmark\,(5831) & \checkmark\,(6559) & \checkmark\,(6890) \\
$T_{2,11}$ & torus & \checkmark\,(22) & \checkmark\,(733) & \checkmark\,(813) & \checkmark\,(5667) & \checkmark\,(6260) & \checkmark\,(6519) \\
$T_{3,7}$ & torus & \checkmark\,(53) & \checkmark\,(1797) & \checkmark\,(1953) & \checkmark\,(13600) & \checkmark\,(14696) & \checkmark\,(15199) \\
$T_{4,5}$ & torus & \checkmark\,(68) & \checkmark\,(2356) & \checkmark\,(2558) & \checkmark\,(18729) & \checkmark\,(20195) & \checkmark\,(20883) \\
K1 & alg & \checkmark\,(104) & \checkmark\,(4141) & \checkmark\,(4414) & \checkmark\,(32732) & \checkmark\,(34655) & \checkmark\,(35558) \\
K3 & alg & \checkmark\,(122) & \checkmark\,(5178) & \checkmark\,(5483) & \checkmark\,(41138) & \checkmark\,(43289) & \checkmark\,(44300) \\
K6a & alg & \checkmark\,(275) & \checkmark\,(12784) & \checkmark\,(13267) & \checkmark\,(100115) & \checkmark\,(103458) & \checkmark\,(105035) \\
K9 & alg & \checkmark\,(564) & \checkmark\,(28584) & \checkmark\,(29309) & \checkmark\,(224004) & \checkmark\,(228995) & \checkmark\,(231354) \\
K8 & alg & \checkmark\,(617) & \checkmark\,(31499) & \checkmark\,(32427) & \checkmark\,(251627) & \checkmark\,(258022) & \checkmark\,(261098) \\
K7 & alg & \checkmark\,(834) & \checkmark\,(43533) & \checkmark\,(44631) & \checkmark\,(343739) & \checkmark\,(351218) & \checkmark\,(354818) \\
K10 & alg & \checkmark\,(4119) & \checkmark\,(227727) & \checkmark\,(230642) & \checkmark\,(1805429) & \checkmark\,(1825082) & \checkmark\,(1834640) \\
$4_1$ & kn & \checkmark\,(10) & \checkmark\,(42) & \checkmark\,(42) & \checkmark\,(223) & \checkmark\,(264) & \checkmark\,(223) \\
\hline
\end{tabular}
}
\end{table}

$T_{2,3}$，$T_{2,5}$ と $\mathbf{4_1}$ の結果は，KN の表の再現である．
KN の表に無いトーラス結び目 5 個（$T_{3,4}$ の 3 箱，$T_{3,5}$，$T_{2,11}$，$T_{3,7}$，$T_{4,5}$）は新しい表である．
ただし KN はトーラス結び目全般について予想しているので，これは予想の確認に当たる．
代数的な iterated torus 結び目 7 個の 3 箱は，本論文の調べた範囲では文献に無い．
3 本ケーブル $K7$ と 2 段の反復 $K10$ では，$|\lambda|=3$ の refined 不変量そのものも文献に見当たらない．
Cherednik--Danilenko~\cite{Cherednik14} は，$K7$ の $\yng(1)$ と $K10$ の DAHA-Betti 多項式を与えている．

\begin{remark}[正規化への依存]
DAHA と楕円 Hall 代数は値を単項式倍を除いてしか決めないので，KN のデータに無い結び目（$T_{3,4}$ の 3 箱，$T_{3,5}$，$T_{2,11}$，$T_{3,7}$，$T_{4,5}$ と全てのケーブル）の判定は，正規化 (N2$'$) に依存する．
依存するのは，$t=q$ で決まらない表現ごとの $(q/t)^{c_\lambda}$ の一点だけである（仮説~\ref{hyp:N2}）．
基準の framing を偶数ずらすことは $x$ のずれに等しく（命題~\ref{prop:N2-framing}），その範囲は表~\ref{tab:framing} の判定に含まれる．
これと異なる割り当ては，$K5c$ の例で Laurent 性と転置対称性を壊す．
\end{remark}

% [R12]
\begin{remark}
入力の $\rCS_{\yng(2,1)}$ 自体は負の係数をもつが，$\hf_{\yng(2,1)}$ は正になる．
この現象はトーラス結び目について KN~\cite{Kameyama17} が指摘したもので（Cherednik~\cite{Cherednik11} の例も参照），本論文の結果はそれが代数的な iterated torus 結び目でも成り立つことを示す．
3 箱の正値性は入力の係数の正値性から直接は従わず，KN の抽出（plethystic な対数と $M$ 行列）を経て初めて現れる．
\end{remark}

\subsection{framing を変えた 3 箱}
% [R13]
\begin{compresult}\label{cr:three-x}
表~\ref{tab:three} の全ての結び目について，3 箱を $x\neq0$ でも判定した．
トーラス結び目 $T_{3,5},T_{2,11},T_{3,7},T_{4,5}$ と代数的なケーブル $K6a,K7,K8,K9$ は $|x|\le10$，$K10$ は $|x|\le2$，$K1$ と $K3$ は $|x|\le2$（$K1$ は $x=\pm6$ も），
自明な結び目と $T_{2,3}$ は $|x|\le6$，$T_{2,5}$ と $\mathbf{4_1}$ は $|x|\le4$，$T_{2,7}$ と $T_{3,4}$ は $|x|\le3$ の範囲である．
いずれも正である（計算結果~\ref{cr:framing-batch} の一部）．
$x$ を正に大きくすると項数が増え，負では減る（$K10$ の $\rho=\yng(3)$ で，$x=-2$ では約 175 万項，$x=2$ では約 186 万項）．
\end{compresult}
3 箱でも $\bt^{-3x}\hf^{(x)}_\rho$ は 7 根の指数和である（命題~\ref{prop:expsum}）．
しかし，6 根のニュートン形では一様な符号の並びが見つからず，全ての $x$ の証明書は得られていない．

\subsection{予想 B とその動機}\label{sec:conjB}
% [R14]
\begin{tcolorbox}
\begin{conjecture}\label{conj:B}
代数的な iterated torus 結び目は，正規化 (N1)(N2$'$) の下の framing 0 で，全ての $\rho$ について正である．
\end{conjecture}
\end{tcolorbox}

KN の物理的な議論は，追加の $U(1)_R$ 対称性を Seifert 多様体の構造から得る．
代数的な iterated torus 結び目の補空間は一般に Seifert 多様体ではないので，この議論は直接は当てはまらない．
予想~\ref{conj:B} の動機は，次の Cherednik--Danilenko~\cite{Cherednik14} の予想にある．
\begin{itemize}
\item 正値性の予想：代数的な結び目と任意の長方形のヤング図形について，DAHA superpolynomial の係数は全て正である（このような結び目を pseudo-algebraic と呼ぶ）．
\item Connection Conjecture (ii)：色付きでない pseudo-algebraic な結び目では，DAHA superpolynomial は（tilde 正規化した）reduced stable Khovanov--Rozansky 多項式に一致する．
\end{itemize}
これらと本論文の結果の関係は次のとおりである．
\begin{itemize}
\item $\yng(1)$：Connection Conjecture (ii) の下で，DAHA の $\yng(1)$ は KhR の Poincaré 多項式 $P_{\mathrm{KhR}}$ の単項式倍 $mP_{\mathrm{KhR}}$ になる．
(N1) により両者は $\bt=-1$ で同じ reduced HOMFLY-PT 多項式を与えるので，$m=\pm\bt^k$ で $\pm(-1)^k=1$ である．両者とも係数が非負なので $m=\bt^{2l}$ である．
$\bt^{2l}$ は $\sigma$ で不変なので，命題~\ref{prop:one} が適用でき，$\yng(1)$ の正値性が従う．
\item 長方形：Cherednik--Danilenko は DAHA の値の正値性を予想し，ホモロジーとの一致を期待している．予想~\ref{conj:B} の長方形の部分は，これを KN の BPS の段に移したものである．
\item $\yng(2,1)$：Cherednik--Danilenko は長方形でない表現の正値性を主張していない．実際 $\rCS_{\yng(2,1)}$ は負の係数をもつ．BPS の段での正値性は，その範囲の外の現象である．
\item 非代数的なケーブル：Cherednik--Danilenko は，正値性をもたない非代数的な iterated torus 結び目があることを認めている．本論文の非代数的なケーブルの多くで，DAHA の値の $\yng(1)$ から負の係数が現れる（付録~\ref{app:nonalg}）．
\end{itemize}

\subsection{双曲結び目の 3 箱へ：非長方形の表現の問題}\label{sec:21}
双曲結び目では，定義に基づく refined $\yng(2,1)$ が文献にもソフトウェアにも無い．
そこで，構造的な条件から $\yng(2,1)$ を決められるかを調べた．
% [R15, R16, R17]
\begin{itemize}
\item 結び目ごと：条件は $t=q$ での 3 本ケーブルから得た $\yng(2,1)$ で色付けた HOMFLY-PT 多項式，ミラー，指数増大性，$U(2)$ への簡約，3 箱の Laurent 性，genus 展開である．
これらを課しても $\yng(2,1)$ は一意に決まらない．解の差は「結び目によらない因子 × 自由な多項式」の形をしている．
曖昧さがあること自体は Anokhina ら~\cite{Anokhina12} も指摘している．
\item 弱い判定：正値性を満たす整数の $\yng(2,1)$ は，$\mathbf{4_1}$（KN の候補）と $\mathbf{6_1}$（813 項の解）に存在する．$\mathbf{5_2}$ では有理数の解しか見つからない（整数解の有無は未決）．
これは予想の検証ではなく，予想と矛盾しないことの確認にとどまる．
\item 族として：$t=q$ ではダブルツイスト結び目の $\yng(2,1)$ について Morozov の式~\cite{Morozov16,Morozov19} がケーブルからの値と一致する．これは既知の式の検証である．
refined の長方形は KNTZ~\cite{Kameyama19} の式を実装し，GNSSS の式と三葉結び目の楕円 Hall 代数の値（$\yng(2,2)$ まで）と照合した．しかし refined の $\yng(2,1)$ は，族の構造の仮定を置いても一意に決まらない．
\end{itemize}
双曲結び目の 3 箱には，定義に基づく $\yng(2,1)$ が必要である．その候補は \S\ref{sec:discussion} で述べる．

\section{検証 IV：4 箱以上}\label{sec:higher}

KN の表は 3 箱までである．この節では，定義どおりの refined 不変量が得られるトーラス結び目について，$|\rho|=4$ から最大 $9$ までの全ての $\rho$ を判定する．
あわせて，$\mathbf{4_1}$ の 4 箱について，未知の $\rCS_{\yng(3,1)}$ に条件を課した弱い判定を述べる．
いずれも framing 0 である．

\subsection{\texorpdfstring{任意の箱数での $\hf_\rho$ の計算}{任意の箱数での hat f の計算}}
KN の $f^q_\mu$ と $\hf_\rho$ の関係 \eqref{eq:f-to-hatf} の行列 $M(t)$ は，$|\mu|\le3$ では明示的に書かれている．
任意の箱数では，次の対角化を用いる．
% [M1]
\begin{lemma}\label{lem:M-diag}
$|\mu|=|\rho|=n$ のとき，$M_{\mu\rho}(t)=\sum_\sigma C_{\mu\sigma\rho}B_\sigma(t)$ は，$B(t)=\sum_\sigma B_\sigma(t)s_\sigma$ との Kronecker 積の行列である．
よって対称群の指標の基底で対角になり，
\[
\sum_\mu M_{\mu\rho}(t)\,\chi_\mu(\kappa)=B(\kappa)\,\chi_\rho(\kappa),\qquad
B(\kappa)=t^{-\frac{n-1}2}\,\frac{\prod_i(1-t^{\kappa_i})}{1-t}
\]
が成り立つ．$\kappa$ は共役類（$n$ の分割）である．
\end{lemma}
\begin{proof}
$C_{\mu\sigma\rho}=\sum_\nu z_\nu^{-1}\chi_\mu(\nu)\chi_\sigma(\nu)\chi_\rho(\nu)$ と指標の直交性から，左辺は $\chi_\rho(\kappa)\sum_\sigma B_\sigma(t)\chi_\sigma(\kappa)$ に等しい．
$B_\sigma$ は hook $\sigma=(n-d,1^d)$ でだけ 0 でなく $(-t)^d t^{-(n-1)/2}$ であり，
$\sum_{d}(-u)^d s_{(n-d,1^d)}=\sum_\kappa z_\kappa^{-1}p_\kappa\prod_i(1-u^{\kappa_i})/(1-u)$ である．
\end{proof}
したがって $f^q$ の冪和基底での係数を $B(\kappa)$ で割り，指標で Schur 基底に戻せば $\hf_\rho$ が得られる．
plethystic な対数は，冪和基底での次数ごとの漸化式で評価点ごとに計算する．
$\rCS_\lambda$ は楕円 Hall 代数の値に計算結果~\ref{cr:M} の閉じた正規化 $M_\lambda=M_\square^{|\lambda|}T_\lambda^{mn}$ を掛けて作る．
この手順は，保存済みの 3 箱の $\hf_\rho$（$T_{2,3}$，$T_{3,4}$ の全 6 つの $\rho$）を全係数で再現し，$f^{\bar t}=M(q^{-1})\hf$ は全ての評価点で成り立った．
計算の詳細は付録~\ref{app:computation} に述べる．

\subsection{トーラス結び目}
% [R18]
\begin{compresult}\label{cr:higher}
表~\ref{tab:higher} のトーラス結び目 8 個と箱数で，全ての $\rho$ について $\hf_\rho$ は正である（framing 0）．
すなわち，同じ全体の符号で spin 表示 $\hbN_{\rho;i,j,k}$ と genus 表示 $\hbN_{\rho,g,\beta,F}$ がともに非負整数になる．
判定した $(K,\rho)$ の組は，$|\rho|\ge4$ で計 340 個である．
\end{compresult}

\begin{table}[h]
\centering
\caption{4 箱以上の判定（framing 0）．「件数」は判定した $\rho$ の数の合計．項数は箱数が最大のときの spin 表示の項数の和．
いずれも spin 表示と genus 表示の両方で負の項は 0 である．}\label{tab:higher}
\begin{tabular}{lcrrr}
\toprule
結び目 & 箱数 $|\rho|$ & 件数 & 最大の箱数での項数 & 最大の箱数の計算時間（秒） \\
\midrule
$T_{2,3}$ & $4$--$9$ & 90 & 1786668 & 428 \\
$T_{2,5}$ & $4$--$8$ & 60 & 2880410 & 299 \\
$T_{2,7}$ & $4$--$8$ & 60 & 6556642 & 1260 \\
$T_{2,9}$ & $4$--$8$ & 60 & 11740842 & 2036 \\
$T_{3,4}$ & $4$--$6$ & 23 & 1358553 & 280 \\
$T_{3,5}$ & $4$--$6$ & 23 & 2336943 & 322 \\
$T_{3,7}$ & $4$--$5$ & 12 & 1307758 & 64 \\
$T_{4,5}$ & $4$--$5$ & 12 & 1848265 & 474 \\
\bottomrule
\end{tabular}
\end{table}

箱数の範囲は次の考えで選んだ．
4 箱で，hook でない $\rho=\yng(2,2)$ と，$\yng(2),\yng(1,1)$ の Adams 作用（$(a,q,t)\mapsto(a^2,q^2,t^2)$）が初めて $\hf_\rho$ に入る．
6 箱で，$d=2,3,6$ の Adams 作用が 3 箱以下の全ての表現を通して出そろう．
7 箱以上では新しい種類の項は現れず，同じ構造の規模が大きくなる．
結び目はブレイドの本数（$a$ 次数の幅）で選んだ．$T_{2,2k+1}$ は thin な族，$T_{3,n}$ は thick で $a$ 次数の幅が 2，$T_{4,5}$ は幅が 3 である．

\begin{remark}
全体の符号は，$|\rho|\equiv2\pmod 4$ で $-1$，それ以外で $+1$ になった（同じ箱数の全ての $\rho$ で共通）．
$|\rho|$ が奇数のときの $a$ 次数の半整数のずらしは，KN の $|\rho|=1,3$ の規約を周期的に延長した（$\beta$ を一様にずらすだけで，正値性には影響しない）．
genus 表示の係数は箱数とともに急速に大きくなる．$T_{2,3}$ の 9 箱では最大 $10^{30}$ 程度である．
3 箱の $T_{4,5}$ でも $10^{21}$ 程度であり，同じ傾向である．
\end{remark}

\subsection{\texorpdfstring{$\mathbf{4_1}$ の 4 箱：弱い判定}{4\_1 の 4 箱：弱い判定}}\label{sec:fig8-four}
$\mathbf{4_1}$ の 4 箱の入力のうち，長方形の $\yng(4),\yng(2,2),\yng(1,1,1,1)$ は KNTZ の refined の式~\cite{Kameyama19} から得られる．
$\yng(1)$ から $\yng(1,1,1)$ では KN のノートブックの値と，対称表現 $[r]$（$r\le4$）では GNSSS の式と全係数一致する．
$\yng(2,1)$ には KN の候補を用いる．
$\yng(3,1)$（と転置で決まる $\yng(2,1,1)$）の refined の値は文献に無い．
そこで $\rCS_{\yng(3,1)}$ の係数を未知数とし，次の一次の条件を課す．
\begin{itemize}
\item $t=q$ で Morozov~\cite{Morozov16} の $\yng(3,1)$ の微分展開に一致する（$\mathbf{4_1}$ では $Z$ 因子の和）．
\item 転置 $\rCS_{\yng(2,1,1)}(a,q,t)=\rCS_{\yng(3,1)}(a,t^{-1},q^{-1})$，両手型性．
\item 指数増大性：$q=1$ で $\rCS_{\yng(1,1)}\rCS_{\yng(1)}^2$，$t=1$ で $\rCS_{\yng(3)}\rCS_{\yng(1)}$ に一致する．
\item $a=t^2$ で $\rCS_{\yng(3,1)}=\rCS_{\yng(2)}\rCS_{\yng(1,1)}$（$U(2)$ への簡約．仮説）．
\item $|\rho|=4$ の全ての $\hf_\rho$ が Laurent 多項式で，genus 展開が閉じ，正である．
\end{itemize}
$\hf_\rho$ は $\rCS_{\yng(3,1)}$ について一次なので（分母 $D$ は $q,t$ だけの多項式），これらは線形計画の条件になる．
% [R19]
\begin{compresult}\label{cr:fig8-four}
上の条件（未知数の範囲は特殊化の値から決めた 2367 個の単項式）の下で，次が成り立つ．
\begin{itemize}
\item 全体の符号 $\sigma=-1$ では，genus 展開の前の正値性だけで実行不可能である．
\item $\sigma=+1$ では，genus 展開の前後の正値性を満たす有理数の解が存在する．
\item 整数の解は，混合整数計画を 50 分実行しても見つからなかった（未決）．
\end{itemize}
\end{compresult}
これは予想の検証ではなく，4 箱でも予想と矛盾しない $\rCS_{\yng(3,1)}$ が（有理数の範囲で）存在することの確認にとどまる．
結論は KN の $\yng(2,1)$ の候補，$U(2)$ の仮説，未知数の範囲に依存する．
定式化は，$\rCS_{\yng(3,1)}$ が楕円 Hall 代数から分かる $T_{2,3}$ で検査した．真の値は全ての条件と未知数の範囲を満たし，そこから組み立てた $\hf_\rho$ は直接の計算と全係数一致した．

\subsection{\texorpdfstring{KN の $[2,1]$ の候補の構造}{KN の [2,1] の候補の構造}}\label{sec:fig8-21}
KN は $\mathbf{4_1}$ の $\yng(2,1)$ の候補を，ミラー対称性，指数増大性，両手型性を課して求めたと述べている~\cite{Kameyama17}．
その候補は，Anokhina らの $Z$ 因子展開~\cite{Anokhina12} の係数を一つ取り替えたものとして書ける．
$\mathfrak{Z}_{I|J}^{(s|\sigma)}=\{AQ^{I+s}T^{-\sigma}\}\{AQ^sT^{-J-\sigma}\}$，$\{x\}=x-x^{-1}$ とおく．
% [R20]
\begin{compresult}\label{cr:fig8-21}
KN の候補は
\[
1+\bigl(\mathfrak Z^{(-1|-1)}_{3|3}+\mathfrak Z^{(1|0)}_{1|1}+\mathfrak Z^{(0|1)}_{1|1}\bigr)
+\bigl(\mathfrak Z^{(0|-1)}_{3|3}\mathfrak Z^{(1|0)}_{1|1}+\mathfrak Z^{(-1|0)}_{3|3}\mathfrak Z^{(0|1)}_{1|1}+\alpha\,\mathfrak Z^{(1|0)}_{1|1}\mathfrak Z^{(0|1)}_{1|1}\bigr)
+\mathfrak Z_{3|3}\mathfrak Z^{(1|0)}_{1|1}\mathfrak Z^{(0|1)}_{1|1}
\]
で $\alpha=1-\{Q\}\{T\}$ としたものに等しい．
Anokhina らは同じ式で $\alpha=1-(Q-T^{-1})(T-Q^{-1})$ とし，係数が全て正の superpolynomial を与えた．
KN の候補は太字の変数で 69 項中 10 項が負で，$\ba=\bq=\bt=1$ での値は $125=5^3$（$\yng(1)$ の値の 3 乗）である．Anokhina らの式では 189 である．
\end{compresult}
$t=q$ では二つの $\alpha$ はともに $1-(Q-Q^{-1})^2$ になる．
つまり KN の候補は，superpolynomial の正値性の代わりに指数増大性を課した選び方に当たる．

さらに，Anokhina らが hook の表現 $\lambda=[r,1^{s-1}]$ に提案した Heegaard--Floer 型の条件
\begin{equation}\label{eq:HF}
\rCS_\lambda\bigl(A=T/Q\,\big|\,Q,T\bigr)=\rCS_{\yng(1)}\bigl(A=T/Q\,\big|\,Q^rT^{s-1},\,T^sQ^{r-1}\bigr)
\end{equation}
を調べた（$A=a^{1/2}$，$Q=q^{1/2}$，$T=t^{1/2}$）．
% [R21]
\begin{compresult}\label{cr:HF}
\eqref{eq:HF} は，$\mathbf{4_1}$ の $\yng(2,1)$（KN の候補と Anokhina らの式の両方）で成り立つ．
楕円 Hall 代数から得た $T_{2,3},T_{2,5},T_{3,4}$ の $\yng(2,1)$ と $\yng(3,1)$ では，単項式 $(T/Q)^{\delta(|\lambda|-1)}$ の違いを除いて成り立つ（$\delta$ は種数型の量）．
\end{compresult}
\eqref{eq:HF} は上の線形計画に含めていない一次の条件である．

KN の方法を $\yng(3,1)$ に延ばすことも試みた．
$t=q$ の Morozov の値を，Anokhina らの形（箱の部分集合ごとに，シフトした $Z$ 因子の積を 1 項置き，係数の補正は箱がつながっていない部分集合だけに付ける）に分解しようとした．
この探索は $\yng(2,1)$ では Anokhina らの展開をただ一つの解として再現する．
しかし $\yng(3,1)$ では，この形の分解は見つからなかった（シフト $|k|\le2$ の整数計画が実行不可能，シフトが箱の組ごとの寄与の和である場合も解なし）．
$\yng(3,1)$ には，部分集合ごとの単純な積を超える構造（Morozov の式の $[6]/[3]$ の係数や異常な成分）が必要と考えられる．

\section{議論と未解決の問題}\label{sec:discussion}

\paragraph{全ての framing での正値性．}
個々の結び目では，$x$ の指数和とニュートン形によって，全ての $x$ での正値性を有限個の判定に帰着できた（定理~\ref{thm:two-all-x}）．
しかしこれは結び目ごとの判定で，一般の結び目についての定理ではない．
ニュートン形の係数の正値性を，入力の性質（ホモロジーの非負性と $\sigma$ 不変性など）から導けるかは未解決である．
ツイスト結び目で $p$ と $x$ が逆符号の領域と，3 箱の全ての $x$ の証明書も得られていない．
また，Cheng--Sułkowski~\cite{Cheng21} の toric brane の refined open BPS 不変量と，framing をつけた自明な結び目の $\hf_\rho$ の対応は確かめていない．

\paragraph{framing の物理的な導出．}
Mariño--Vafa の符号が Laurent 性に必要であることは示した（命題~\ref{prop:sign}）．
しかし，refined の平坦座標のずれ $V\mapsto-(q/t)^{1/2}V$ を幾何から導くことはしていない．
弦側で $U(1)_R$ の作用が framing によらず自明であることも仮定にとどまる．

\paragraph{双曲結び目の非長方形の表現．}
双曲結び目の 3 箱には，定義に基づく refined $\yng(2,1)$ が必要である．候補は三つある．
(i) refined の Racah 行列による構成（arborescent 計算の refined 化）．
(ii) DAHA・楕円 Hall 代数型の定義の双曲結び目への拡張（種数 2 の DAHA は $\mathfrak{sl}_2$ に限られる~\cite{Arthamonov15}）．
(iii) $\yng(2,1)$ で色付けた HOMFLY-PT ホモロジーの計算と，KN の $\overline{\rCS}$ への変換規則．
いずれも本論文の範囲を超える．
$\mathbf{4_1}$ の 4 箱では，$\yng(3,1)$ も同じ問題に当たる（\S\ref{sec:fig8-four}）．
KN の $\yng(2,1)$ の候補は Anokhina らの $Z$ 因子展開に乗るが（計算結果~\ref{cr:fig8-21}），同じ形の展開は $\yng(3,1)$ では見つからなかった．
Heegaard--Floer 型の条件 \eqref{eq:HF} は楕円 Hall 代数の値でも成り立つので，$\yng(3,1)$ を絞る新しい条件として使える．

\paragraph{正規化．}
単項式の正規化の閉じた式（計算結果~\ref{cr:M}）は，refined 衛星公式（仮説~\ref{hyp:sat}）の下で導出した．
DAHA と楕円 Hall 代数は単項式を決めないので，この仮説は物理的な仮説として残る．
トーラス結び目の部分では，Mellit の系の標準の正規化での形を計算で確かめた（計算結果~\ref{cr:mellit}）．
$d_\pm$ の表現の元と Bergeron--Garsia--Leven--Xin の $Q_{m,n}$ の対応を具体的に示すことが残っている．
ただし，これは正値性の結果には影響しない．

\paragraph{箱数．}
トーラス結び目の 4 箱以上の判定は，箱数と結び目を有限の範囲に限った計算である（表~\ref{tab:higher}）．
genus 表示の係数は箱数とともに急速に大きくなるが，負の項は現れなかった．
全ての箱数での正値性を，楕円 Hall 代数の構造から示すことは未解決である．

\paragraph{そのほか．}
ダブルツイスト結び目の異符号の領域の残り（\S\ref{sec:two}）と，2 箱の単峰性の構造的な理由も未解決である．


\appendix
\section{計算法の詳細}\label{app:computation}
\paragraph{評価と補間．}
$\overline{\rCS}_\lambda$ と $\hf_\rho$ は，31 ビットの素数 $p$ を法として格子の点で評価し，$(A,Q,T)$ の各変数について補間して求める．
係数が $p/2$ を超えるときは 2〜3 個の素数で補間し，中国剰余定理（Garner の方法）で復元する．
得られた多項式は，補間に使っていない素数の点で検証する．
評価の点はまとまりごとに保存し，計算を途中から再開できるようにした．
\paragraph{genus 展開．}
補題~\ref{lem:G} と同じく $v=\bq^2\bt$，$w=v+v^{-1}$，$w_0=\bt+\bt^{-1}$ とおくと $G=w-w_0$ である．各行を $w$ の多項式に直し，Horner 法で $G$ の冪に展開する．
\paragraph{4 箱以上．}
$|\rho|\ge4$ では $\hf_\rho$ を評価点ごとに直接計算する（\S\ref{sec:higher}）．
$A$ も数値の点で評価し，$a=A^2$ の多項式として補間する（範囲の外の 1 点で確かめる）．
トーラス結び目では $Q,T$ の指数の偶奇が一様なので，$q=Q^2,t=T^2$ で補間して評価点を 1/4 にした．
楕円 Hall 代数の計算では，冪和 $p_\rho$ に対応する演算子の作用とその評価を表現 $\lambda$ をまたいで共有した．
計算時間は表~\ref{tab:higher} のとおりで，8 プロセスの並列で行った．$T_{2,7}$ と $T_{2,9}$ の 8 箱は 1 件 10 分を目安とした見積もりを超えた（前者は他の計算との競合による）．
\paragraph{計算時間．}
計算は Apple M3 Pro（メモリ 36 GB）と Mac mini M1（メモリ 8 GB）で行った．
最も重いのは $K10$ である．$\rCS_{\yng(3)}$ は 2 つの素数で約 4 時間（14440 秒），$\rCS_{\yng(2,1)}$ は約 2 時間（7387 秒）かかった．
3 箱の $\wh f_\rho$ は 3 つの素数を要し，4050 秒，最大メモリ 17.6 GB であった．
$K7$ と $K8$ の $|\rho|\le3$ は，それぞれ 12275 秒と 22853 秒で，最大メモリは 2 GB 未満であった．
\paragraph{実装の検証．}
実装は，小さな入力で各段を確かめるテストで検証した．
対象は，KN の値の再現，Cherednik--Danilenko の明示例，ブレイドから計算した HOMFLY-PT 多項式，補間と中国剰余定理，genus 展開の二つの実装の一致などである．

\section{正規化の導出}\label{app:normalization}
\paragraph{(A) ケーブル．}
framing 演算子 $\tau:P_\mu\mapsto T_\mu P_\mu$ について，$\tau^jP_\lambda[p_k\mapsto P_{k(m,n)}]\cdot1=P_\lambda[p_k\mapsto P_{k(m,n+jm)}]\cdot1$ が成り立つ（計算結果~\ref{cr:tau}）．
ケーブルの構成は companion について線形である．companion に (N2$'$) を適用すると，パターンの傾きが $(m,-s')$ から $(m,-n)$（Seifert framing）にずれる．
仮説~\ref{hyp:sat} と合わせると，$M_{C_{m,n}K}=M_K^mM_{T_{m,n}}$ と $F=mn$ が従う．
\paragraph{(B) トーラス結び目．}
(I) Gorsky--Neguţ の正規化前の reduced 値の $a^0$ の係数は $(-t/q)^{m-1}C_{m,n}(1/q,t)$ である（$C_{m,n}$ は有理 $q,t$-Catalan 数）．
これは Gorsky--Neguţ の予想であり，Bergeron--Garsia--Leven--Xin の $Q_{m,n}$ を介して Mellit~\cite{Mellit16} の定理から従う．
(II) 最低 $a$ 次数 $\delta$ は Franks--Williams の正のブレイドについての定理から従う．
(III) $t=q$ での符号と $q$ の冪は，$C_{m,n}$ の $q,t$ 対称性から決まる．
(IV) $(q/t)$ の指数は，Mellit~\cite{Mellit17} の Dyck 路の表示と計算結果~\ref{cr:mellit} の式から決まる．
以上から $M_{T_{m,n}}=(-q/t)^{m-1}(a(q/t)^{1/2})^{\delta}$ である．

\section{非代数的なケーブル}\label{app:nonalg}
本論文で扱う非代数的なケーブルは，$T_{2,3}$ の $(2,1)$-，$(2,3)$-，$(2,5)$-，$(2,7)$-，$(2,-1)$-ケーブル（順に $K2,K4,K5a,K5b,K5c$）と，$T_{2,5}$ の $(2,1)$-ケーブル（$K6b$）である．
非代数的なケーブル（$n_2<m_1n_1m_2$）では，DAHA の値はホモロジーの Poincaré 多項式ではない（スペクトル系列の必要条件を満たさない）．
したがって KN の正値性の前提を満たさない．表~\ref{tab:nonalg} に判定を示す．
\begin{table}[h]
\centering
\caption{非代数的なケーブルの判定（framing 0）．各欄は「spin 表示の負の項の数/項数; genus 表示の負の項の数/項数」（全体の符号はそれぞれで良い方）．}\label{tab:nonalg}
\resizebox{\textwidth}{!}{% scripts/framed_knot/paper_summary.py が生成。\checkmark(n) は spin 表示と genus 表示の両方で負の項 0（n は spin 表示の項数）。それ以外は spin 表示の負/項数; genus 表示の負/項数（前後それぞれで全体の符号は良い方）。framing 0。
\begin{tabular}{llcccccc}
\hline
結び目 & 分類 & $\yng(1)$ & $\yng(2)$ & $\yng(1,1)$ & $\yng(3)$ & $\yng(2,1)$ & $\yng(1,1,1)$ \\
\hline
K2 & non & 10/47; 6/48 & 152/1136; 41/1045 & 174/1273; 51/1190 & 967/8525; 253/7852 & 1078/9527; 292/8774 & 1156/9990; 318/9254 \\
K4 & non & 18/42; 20/43 & 286/1109; 195/952 & 330/1255; 224/1090 & 2561/8614; 1749/7420 & 2850/9646; 1989/8325 & 3025/10121; 2128/8795 \\
K5a & non & 8/38; 5/47 & 409/1226; 332/1066 & 460/1377; 376/1210 & 4441/9758; 3631/8278 & 4982/10844; 4066/9235 & 5164/11348; 4263/9733 \\
K5b & non & \checkmark\,(48) & 136/1395; 66/1305 & 149/1537; 78/1452 & 2097/10706; 1060/9612 & 2358/11707; 1181/10553 & 2351/12202; 1228/11040 \\
K5c & non & 5/57; 0/69 & 78/1307; 0/1299 & 77/1443; 0/1433 & 400/9501; 6/9123 & 423/10494; 7/10053 & 376/10864; 0/10444 \\
K6b & non & 30/133; 10/178 & 576/4400; 81/4367 & 618/4677; 91/4656 & 3626/33214; 512/32175 & 3766/35136; 541/34032 & 3878/36067; 560/34980 \\
\hline
\end{tabular}
}
\end{table}
破れ方は三つの型に分かれる．
(I) $T_{2,3}$ の $(2,3)$，$(2,5)$-ケーブル（$K4$，$K5a$）では，色付きでない HOMFLY-PT の段階で，偶奇のそろった refinement による正値性が不可能である．
(II) $(2,1)$-ケーブル（$K2$，$K6b$）と $(2,-1)$-ケーブル（$K5c$）では，HOMFLY-PT は正値性と両立するが，DAHA の $\rCS_\square$ 自体が負の係数をもつ．
(III) $T_{2,3}$ の $(2,7)$-ケーブル（$K5b$）の色付きの場合の破れは framing に依存する．偶数の framing $x$ を調べると，2 箱は $x\le-12$，3 箱は $x\le-14$ で前後とも正になる（$|x|\le24$ の偶数で計算）．
$K5c$ は展開の後だけが（ほぼ）正になる例で，補題~\ref{lem:G} の逆が成り立たないことを示す．

\section{データとコード}\label{app:data}
計算・判定のコードとテストは \url{https://github.com/kameyama/refined_lmov} で公開している．
このリポジトリには，$\overline{\rCS}_\lambda$，$\hf_\rho$，spin 表示と genus 表示の値も置いてある．ただし 4 箱以上の値は大きいので置いておらず，コードで再計算できる．


\printbibliography
\end{document}
